\providecommand{\toplevelprefix}{../..}  %
\documentclass[../../book-main_zh.tex]{subfiles}

\begin{document}

\chapter{基于有损压缩的表征学习}
\label{ch:lossy-compression}

\begin{quote}
    \hfill ``\textit{我们为学习而压缩，也为压缩而学习}。''
    
    $~$\hfill —— \textit{高维数据分析}, Wright and Ma, 2022
\end{quote}
\vspace{5mm}

让我们回顾一下目前为止讨论的内容。我们已经表明，通过\textit{迭代去噪}将纯噪声（高熵）的分布逐渐转化为要学习的分布，可以刻画并采样具有任意低维支撑的分布，如\Cref{fig:representation-learning}（左）所示。我们已经证明，每次迭代中的去噪器通过Tweedie公式 \eqref{eq:tweedie} 与分布的对数密度梯度直接相关。我们展示了在实践中，这样的\textit{去噪器} \(\bar{\vx}_{\theta}\) 可以使用有限样本进行估计。因此，我们开发了一种学习或追求数据分布的通用方法。

然而，在这种方法中，分布的编码隐含在去噪器的函数形式和参数（如果有的话）中。事实上，敏锐的读者可能已经注意到，对于一般分布，我们从未明确指定去噪器的函数形式应该是什么。在实践中，人们通常使用采用经验性架构设计的深度神经网络对其进行建模。此外，虽然我们知道上述去噪过程会降低熵，但我们不知道降低了多少，也不知道中间分布和最终分布的熵。因此，了解去噪操作的（最优）形式以及如何测量熵的变化将非常有用。最后，我们在上一章末尾表明，如果去噪过程将熵降低得\textit{太多}，则会导致仅学习训练数据，而无法进行任何插值或外推。解决这些问题的方法将在本章中逐步展开。

\paragraph{信息度量。}
回顾一下，我们的总体目标是对来自具有低维支撑的连续分布的数据进行建模。如果我们想找出生成数据的最简单模型，我们可能会考虑三种典型的简约性度量：维度、体积或微分熵。如果我们使用维度，那么显然给定数据集的最佳模型是经验分布本身，它是零维的。\footnote{在引言\Cref{ch:intro}中，我们提出智能的一个重要目标是追求观测数据中的低维结构。然而在技术层面上，我们不应从字面上理解“低维”一词。本章的目的正是要准确阐明我们应该追求什么样的低维结构以及如何追求。}如果我们改用数据的熵，那么对于所有具有低维支撑的分布，微分熵总是负无穷大。或者如果我们使用数据所张成空间的体积，任何低维模型支撑的体积总是零。那么，在所有可能生成相同数据样本的具有低维支撑的分布中，我们如何基于这些根本无法相互区分的简约性度量来决定哪一个更好呢？

在本章中，我们将讨论一个框架，该框架通过将学习到的分布与显式可计算的编码和解码方案相关联，来缓解上述技术困难，这遵循并进一步明确了\Cref{sub:min_entropy}末尾建议的第二种方法。正如我们将看到的，这种方法本质上允许我们根据与编码方案相关的\textit{有损编码率}来准确近似学习分布的熵，正如在\Cref{sec:diffusion_memorization}末尾所描述和启发（但未形式化）的那样。粗略地说，这种度量可以被视为经验数据分布体积的广义概念。有了这样的度量，我们就可以准确地量化通过分布的变换（例如通过去噪器）降低了多少熵。此外，有了这样的度量，我们至少在原则上可以推导出最快降低熵的最优去噪算子。

\paragraph{表征学习。}
尽管上一章介绍的扩散和去噪过程在原则上允许我们学习任意的低维分布，但它并没有给出分布的显式表征。正如我们将看到的，学习到的分布将被反复访问并用于后续任务（例如用作\Cref{ch:conditional-inference}中贝叶斯推断的先验）。因此，我们希望能够通过映射 $f(\vx,\theta)$ 将数据分布转换为某种便于高效访问的标准结构化形式，例如（非相干）线性子空间或低秩高斯分布的混合，如\Cref{fig:representation-learning}右侧所示。我们将在本章中将此过程形式化为“表征学习”。

\begin{figure}
    \centering
    \includegraphics[width=0.8\linewidth]{chapters/chapter4/figs/manifold-linearization.pdf}
    \caption{为数据分布（左）学习结构化表征（右）。请注意，除了像在上一章\Cref{ch:denoising}中那样学习（和采样）原始数据的分布之外，我们在这里还希望学习数据分布的变换 $f(\vx, \theta)$，从而产生具有所需更好几何或统计特性（例如线性和非相干性）的表征。}\label{fig:representation-learning}
\end{figure}

正如我们还将看到的，对于任意分布，上述信息度量（编码率）并不（容易）计算。但对于子空间或低秩高斯分布的混合，我们可以推导出该度量的精确闭式解。通过优化该度量，我们将能够找到最优表征。正如我们将在下一章\Cref{ch:deep-networks}中看到的，优化该度量还允许我们推导出以最有效方式执行变换 $f(\vx,\theta)$ 的算子的显式形式。这将为深度网络的架构提供原则性的解释，并带来更高效的深度架构设计。

\section{通过有损编码进行压缩} \label{sec:lossy_compression}

在本节中，我们将\textit{有损编码率}发展为一种可计算的复杂度度量，用于支撑为低维的分布——在这种情形下，熵、维度和体积等经典度量都无法区分候选模型。我们首先针对这些病态现象说明有损编码的必要性（\Cref{sub:rate-dist}），随后展示失真编码率函数如何通过球覆盖将该编码论度量与数据支撑的几何结构联系起来（\Cref{thm:covering-number-rate-distortion}）。这一联系是连接信息论方法与几何方法、用于学习低维分布的桥梁。最后，我们为高斯分布给出有损编码率的紧致闭式近似（\Cref{subsec:lossy DR}），并给出用于低维高斯混合聚类的有效算法（\Cref{sec:clustering-Gaussians}），该算法将在后续小节中推广到非线性分布。

\subsection{有损编码的必要性}\label{sub:rate-dist}

我们之前曾多次讨论过一个困难：如果我们最终从有限样本中学习分布，并且我们的去噪器函数类包含足够多的函数，我们如何确保我们从\textit{真实}分布（具有低维支撑）中采样，而不是从可能以高概率产生这些有限样本的任何其他分布中采样？让我们通过一些具体的例子来揭示一些概念和技术上的困难。

\begin{example}[体积、维度与熵]\label{eg:measures-of-complexity}
    对于\Cref{fig:1d-line}顶部所示的例子，假设我们从一条直线（例如在二维平面中）上的均匀分布中提取了一些样本。直线或样本集的体积为零。在几何上，生成的有限样本集上的经验分布是能够产生该有限样本集的{\em 最小维度}分布。\footnote{一组离散样本都是零维的，而支撑直线是一维的。}但这似乎与另一种复杂性度量——熵——形成了对比。直线的（微分）熵是负无穷大，但这个样本集的（离散）熵是有限且为正的。因此，根据这些常见的简约性或复杂性度量，我们似乎面临一种自相矛盾的情况：它们根本无法正确区分具有低维支撑的分布（的模型），有些甚至似乎以完全相反的方式区分它们。\footnote{当然，严格来说，微分熵和离散熵是不能直接比较的。}
\end{example}

\begin{example}[密度]\label{eg:density} 考虑\Cref{fig:1d-line}中所示的两组采样数据点。在几何上，它们本质上是相同的：每组由八个点组成，每个点以相等的频率 $1/8$ 出现。唯一的区别是，对于第二个数据集，一些点足够“接近”，可以被视为在它们各自的“簇”周围具有更高的密度。哪一个与可能生成这些样本的真实分布更相关？我们如何调和在解释这种（经验）分布时的这种歧义？
    \begin{figure}[t]
        \centering
        \includegraphics[width=0.7\linewidth]{\toplevelprefix/chapters/chapter4/figs/one-dim-distribution.png}
        \caption{在直线上观察到的具有不同几何构型的八个点；均匀分布（上）或紧密聚集（下）。}
        \label{fig:1d-line}
    \end{figure}
\end{example}

构建数据集的显式编码和解码方案还存在另一个技术困难。给定采样数据集 $\vX = [\vx_1, \ldots, \vx_N]$，如何设计一种可在具有有限内存和计算资源的机器上实现的编码方案？请注意，即使表示一个一般的实数也需要无限数量的数字或比特。因此，人们可能会想，分布的熵是否是其（最优）编码方案复杂性的直接度量。我们用另一个简单的例子来检验这个问题。
\begin{example}[精度] \label{eg:two-inrational}
    考虑一个离散分布 $\vX = [e, \pi]$，以 $1/2$ 的等概率取欧拉数 $e \approx 2.71828$ 或圆周率 $\pi \approx 3.14159$ 的值。该分布的熵为 $H =1$，这表明人们可以分别用一位数字 $0$ 或 $1$ 来编码这两个数字。但是你能在有限状态机上实现该代码的解码方案吗？答案实际上是否定的，因为精确描述这两个数字中的任何一个都需要无限多的比特。
\end{example}

因此，通常不可能有一种编码和解码方案能够精确地重现来自任意实值分布的样本。\footnote{也就是说，如果想要精确地编码这些样本，唯一的方法就是记住每一个样本。}但是，如果对分布进行编码却无法对从同一分布中抽取的样本进行解码，那么其实际价值微乎其微。

因此，为了确保任何编码/解码方案都是可计算的，并且可以在有限的内存和计算资源下实现，我们需要对样本 $\vx$ 进行量化，并仅在一定精度（例如 $\epsilon > 0$）内对其进行编码。{\em 通过这样做，本质上，如果任意两个数据点之间的距离小于 $\epsilon$，我们就认为它们是等价的。}更准确地说，我们希望考虑编码方案
\begin{equation}
    \vx \mapsto \hat \vx
\end{equation}
使得由量化引起的预期误差以 $\epsilon$ 为界。在数学上更方便且在概念上几乎相同的是，将预期的\textit{平方}误差限制在 \(\epsilon^{2}\) 内，即
\begin{equation}
    \Ex[d(\vx, \hat \vx)^{2}] \le \epsilon^{2}.
\end{equation}
通常，距离 \(d\) 被选择为欧几里得距离，或2-范数。\footnote{更一般地，我们可以用任何所谓的\textit{散度}（divergence）来代替 \(d^{2}\)。}我们在下文中将采用这一选择。

\subsection{失真编码率与数据几何}
当然，在所有满足上述约束的编码方案中，我们希望选择使产生的编码率最小化的方案。对于给定的随机变量 $\vx$ 和精度 $\epsilon$，该编码率称为\textit{失真编码率}（rate distortion），记为 $\cR_{\epsilon}(\vx)$。

为了更好地理解这种失真编码率，我们求助于信息论中的一个深刻定理，该定理最初由 \textcite{shannon1959coding} 证明，它将这种失真编码率与其他信息论量联系起来。为了理解这个定理，我们首先需要定义两个分布之间\textit{差异}的度量，称为\textit{Kullback-Leibler} (KL) 散度，或\textit{相对熵}，记为 \(\KL(\cdot \mmid \cdot)\)，并结合熵和交叉熵（见\Cref{sub:min_entropy}）定义为
\begin{align}
    \KL(p \mmid q) &= \CE(p \mmid q) - h(p) = \Ex_{\vx \sim p}[-\log
    q(\vx)] - \Ex_{\vx \sim p}[-\log p(\vx)] \\
    &= \Ex_{\vx \sim p}[\log (p(\vx)/q(\vx))].
\end{align}
KL散度总是非负的，当且仅当 \(p = q\) 时等于零；\footnote{从技术上讲，这个等式应该被理解为“几乎处处”，即除了可能在一个零测度（体积）的集合上，因为这个集合不会影响任何积分的值。}在 \(p\) 和 \(q\) 具有相同支撑且都是概率密度的情况下，这可以通过以下基于Jensen不等式和函数 \(\log(\cdot)\) 严格凹的事实的计算来证明：
\begin{align}
    - \KL(p\mmid q)
    &= -\Ex_{\vx \sim p}[\log(p(\vx)/q(\vx))] = \Ex_{\vx \sim
    p}[\log(q(\vx)/p(\vx))] \\
    &\leq \log \Ex_{\vx \sim p}[q(\vx)/p(\vx)] =  \log \int_{\R^{D}}
    \frac{q(\vxi)}{p(\vxi)}p(\vxi)\odif{\vxi} \\
    &= \log \int_{\R^{D}}q(\vxi)\odif{\vxi} = \log \Ex_{\vx \sim
    q}[1] = \log 1 = 0.
\end{align}
然后，我们定义\textit{互信息} \(I(\cdot; \cdot)\)，它是两个随机变量之间\textit{依赖性}的度量，如下所示：
\begin{equation}\label{eq:mutual-information}
    I(\vx; \hat{\vx})
    = \KL(p(\vx, \hat{\vx}) \mmid p(\vx) p(\hat{\vx})),
\end{equation}
有了前面的定义，我们可以用纯概率术语将失真编码率 \(\cR_{\epsilon}\) 写为
\begin{equation}
    \cR_{\epsilon}(\vx)
    = \min_{p(\hat \vx \mid \vx): \Ex[\norm{\vx -\hat \vx}_2^{2}] \le
    \epsilon^{2}}
    I(\vx; \hat{\vx}).
    \label{eqn:rate-distortion-general}
\end{equation}

请注意，\eqref{eqn:rate-distortion-general} 中的最小化是在所有满足失真约束 $\Ex[\norm{\vx- \hat \vx}_2^{2}] \le \epsilon^{2}$ 的条件分布 $p(\hat \vx \mid \vx)$ 上进行的。每个这样的条件分布都会诱导一个联合分布 $p(\vx, \hat{\vx}) = p(\hat\vx \mid \vx) p(\vx)$，它决定了互信息 \eqref{eq:mutual-information}。互信息（以及因此的失真编码率）的许多便利性质都隐含在KL散度的相应性质中。从定义中我们知道，$\cR_{\epsilon}(\vx)$ 是关于 $\epsilon$ 的{\em 非增}函数。

\begin{example}[高斯信源的失真编码率（定理 13.3.3, \cite{Cover-Thomas}）]\label{ex:gaussian_rate_distortion}
    考虑一个高斯随机变量 \(\vx \sim \dNorm(0, \sigma^{2})\)。对于这种情况，众所周知，失真编码率函数可以以闭式解计算为：
    \begin{equation}
        \cR_{\epsilon}(x) =
        \casework{\frac{1}{2}\log\left(
            {\sigma^{2}}/{\epsilon^{2}}\right), & 0 \leq
        \epsilon \leq \sigma, \\ 0, & \epsilon > \sigma}.
    \end{equation}
    这以一种特别有趣的方式扩展到了 \(D\) 维的情况。假设 \(\vx \sim \dNorm(\vzero, \vSigma)\)，其中 \(\vSigma \in \PSD(D)\) 具有特征值 \(\lambda_{1} \geq \cdots \geq \lambda_{D} \geq 0\)。那么失真编码率可以通过“注水算法”（water filling）获得，如下所示：
    \begin{equation}
        \cR_{\epsilon}(\vx) = \frac{1}{2} \sum_{i =
        1}^{D}\log\rp{\frac{\lambda_{i}}{\min\{\kappa, \lambda_{i}\}}},
    \end{equation}
    其中“水位” \(\kappa\) 满足 \(\sum_{i = 1}^{D}\min\{\kappa, \lambda_{i}\} = \epsilon^{2}\)。直观地说，方差 \(\lambda_{i}\) 大于 \(\kappa\) 的方向被有效编码，而方差较小的方向被丢弃，从而确保以最小的互信息满足失真约束。有关此闭式解的详细推导，请参阅\Cref{exer:KKT}。
\end{example}

\begin{remark}
    事实证明，在以下意义下，失真编码率是 $\vx$ 熵的一种可实现近似。假设 $\vx$ 和 $\hat{\vx}$ 是连续随机向量，那么互信息可以写成
    \begin{equation}\label{eq:information-entropy}
        I(\vx; \hat\vx) = h(\vx) - h(\vx \mid \hat \vx),
    \end{equation}
    其中 $h(\vx \mid \hat \vx) = -\bE[\log_2 p(\vx \mid \hat\vx)]$ 是给定 $\hat \vx$ 时 $\vx$ 的\textit{条件熵}。因此，当在所有满足约束 $\mathbb{E} [\norm{\vx- \hat \vx}_2^2] \le \epsilon^2$ 的分布中，$\vx$ 的熵与给定 $\hat \vx$ 时 $\vx$ 的条件熵之差最小化时，即可实现最小编码率。

    事实上，不需要假设 $\vx$ 和 $\hat \vx$ 是连续的也能得出上述类型的结论。例如，如果两个随机向量都是离散的，在对离散值随机向量的KL散度进行适当解释后，我们有
    \begin{equation}\label{eq:mutual}
        I(\vx; \hat\vx) = H(\vx) - H(\vx \mid \hat \vx),
    \end{equation}
    其中 $H(\vx)$ 在\Cref{sub:min_entropy}的 \eqref{eqn:entropy-discrete} 中定义。更一般地，可以使用测度论中的高级数学概念来定义任意随机变量 $\vx$ 和 $\hat \vx$ 的互信息（以及因此的失真编码率），包括那些具有相当奇特的低维分布的随机变量；参见 \textcite[\S 8.5]{Cover-Thomas}。
\end{remark}

\begin{remark}
    给定一组数据点 $\vX = [\vx_1,\ldots, \vx_N]$，我们总可以将其解释为定义在这 $N$ 个向量上的均匀离散分布（每个点的概率均为 $1/N$）的样本。这种分布的熵为 $H(\vX) = \log_2 N$。\footnote{再次注意，即使我们可以用这个编码率对这些向量进行编码，我们也无法以任意精度对它们进行解码。}然而，即使 $\vX$ 在其样本上服从均匀分布，如果这些样本并不均匀分散，而是有许多样本紧密聚集在一起，那么通过有损编码方案可实现的编码率 $\cR_{\epsilon}(\vX)$ 也可能显著低于 $H(\vX)$。因此，对于\Cref{fig:1d-line}中所示的第二个分布，在适当选择量化误差 $\epsilon$ 时，可实现的有损编码率可能显著低于将其编码为均匀分布的编码率。\footnote{尽管如此，对于这种离散均匀分布，当 $\epsilon$ 足够小时，我们总是有 $H(\vX) \approx \cR_{\epsilon}(\vX)$。}还要注意，有了失真编码率的概念，例 \ref{eg:two-inrational} 中讨论的困难也消失了：我们可以选择两个分别足够接近这两个无理数的有理数。由此产生的编码方案将具有有限复杂度。
\end{remark}

\begin{example}\label{example:sphere-covering-rate-distortion}
    有时，我们可能会面临相反的情况，即我们想先固定编码率，然后尝试找到使失真最小化的编码方案。例如，假设我们只想对从分布中采样的点使用固定数量的代码，并且我们想知道如何设计代码，以便在编码/解码方案中使平均或最大失真最小化。例如，给定单位正方形上的均匀分布，我们想知道我们能以多高的精度对从该分布中抽取的点进行编码，比如使用 $n$ 个比特。这个问题等价于问：最小半径（即失真）是多少，使得我们可以用 $2^n$ 个该半径的圆盘覆盖单位正方形。图 \ref{fig:seven-circles-packing} 显示了正方形的近似最优覆盖，其中 \(n = 4, 6, 8\)，因此分别有 \(2^{n} = 16, 64, 256\) 个圆盘。请注意，圆盘的最优半径随着圆盘数量 \(2^{n}\) 的增加而减小。
    \begin{figure}
        \centering
        \includegraphics[width=0.3\textwidth]{\toplevelprefix/chapters/chapter4/figs/16.png}
        \hfill
        \includegraphics[width=0.3\textwidth]{\toplevelprefix/chapters/chapter4/figs/64.png}
        \hfill
        \includegraphics[width=0.3\textwidth]{\toplevelprefix/chapters/chapter4/figs/256.png}

        \caption{使用启发式优化算法计算的覆盖正方形的 \(2^{4}\)、\(2^{6}\) 和 \(2^{8}\) 个圆盘的近似最优解，以及相应的半径。}
        \label{fig:seven-circles-packing}
    \end{figure}
\end{example}

\begin{figure}[t]
    \centering
    \includegraphics[width=0.6\textwidth]{\toplevelprefix/chapters/chapter4/figs/continuation.png}
    \caption{\small\textbf{用 \(\epsilon\)-球近似低维分布。} 我们可以看到，随着 \(\epsilon\) 参数减小，\(\epsilon\)-球的并集越来越好地近似真实分布（黑色）的支撑。此外，通过用高斯混合（每个高斯分布的协方差为 \((\epsilon^{2}/D)\vI\)）近似真实分布而得到的关联去噪器（其输入-输出映射由图中箭头表示），也越来越接近真实去噪器。在较大的 \(\epsilon\) 下，此类去噪器几乎不指向真实分布附近；而在较小的 \(\epsilon\) 下，它们可以紧密近似真实去噪器。\Cref{thm:covering-number-rate-distortion} 表明，这种近似刻画了小失真 $\epsilon$ 下的失真编码率函数，统一了编码率最小化和去噪这两条并行路径，用于学习无病态的低维分布。}
    \label{fig:continuation}
\end{figure}

事实证明，对于一般分布 $p(\vx)$，要获得失真编码率函数 \eqref{eqn:rate-distortion-general} 的闭式表达式通常非常困难。然而，正如\Cref{example:sphere-covering-rate-distortion}所示，在一些重要的特殊情形下，分布 $p(\vx)$ 支撑集的\textit{几何结构}可以与失真编码率函数建立联系，从而与失真水平 $\epsilon$ 下的最优编码率建立联系。事实上，这个例子可以推广到如下设置：$p(\vx)$ 的支撑是足够规则的紧集（包括低维分布），且 $p(\vx)$ 在其支撑上均匀分布。这覆盖了大量具有实际意义的情形。我们在下面的结果中将这一思想形式化，并为一个特例确立该性质。

\begin{theorem}\label{thm:covering-number-rate-distortion}
    假设 \(\vx\) 是一个随机变量，使得其支撑 \(K \doteq \Supp(\vx)\) 是一个紧集。将\textit{覆盖数} \(\cN_{\epsilon}(K)\) 定义为能够覆盖 \(K\) 的半径为 \(\epsilon\) 的球的最小数量，即，
    \begin{equation}
        \cN_{\epsilon}(K) \doteq \min\left\{n \in \bN \colon \exists
            \vp_{1}, \dots, \vp_{n} \in K\ \text{s.t.}\ K \subseteq
        \bigcup_{i = 1}^{n}B_{\epsilon}(\vp_{i})\right\},
    \end{equation}
    其中 \(B_{\epsilon}(\vp) = \set{\vxi\in \R^D \given \norm{\vxi -
            \vp}_2 \leq
    \epsilon}\) 是以 \(\vp\) 为中心、半径为 \(\epsilon\) 的欧几里得球。那么有
    \begin{equation}
        \cR_{\epsilon}(\vx)
        \leq \log_{2} \cN_{\epsilon}(K).
    \end{equation}
    如果此外，$\vx$ 在 $K$ 上均匀分布，并且 $K$ 是相互正交的低维超球面的并集，\footnote{事实上，也可以对高度不规则的 $K$（例如紧流形甚至分形）得到类似结果，但其陈述会更技术化：参见 Riegler 等人 \cite{Riegler2018-jh,Riegler2023-rr}。我们在\Cref{app:rate-distortion-covering}中给出一个简单证明，展示了正交超球面并集情形；当每个子空间维度足够高时，这一假设与“支撑正交的低秩高斯混合”是等价的。更多讨论见\Cref{app:rate-distortion-covering}。}那么存在一个匹配的下界：
    \begin{equation}
        \cR_{\epsilon}(\vx)
        \geq
        \log_{2} \cN_{\epsilon}(K) - O(D).
    \end{equation}
\end{theorem}
\begin{proof}
    该定理的证明超出了本书的范围，我们将其推迟到\Cref{app:rate-distortion-covering}。
\end{proof}

\Cref{thm:covering-number-rate-distortion} 的含义可以总结如下：对于 $\vx$ 分布的足够精确的编码，最小失真编码率框架完全由 $\vx$ 支撑上的球堆积问题来表征。\Cref{thm:covering-number-rate-distortion} 证明的核心确实可以推广到更复杂的分布，例如足够非相干的流形混合，但我们将其留待未来的研究。因此，失真编码率可以被认为是一种“概率感知”的方式，通过许多小球的混合来近似 \(\vx\) 分布的支撑。我们现在讨论这与我们在本章前面讨论的去噪-扩散-熵复杂性层次结构之间的另一个联系。

\begin{remark}\label{rem:slb}
    \Cref{thm:covering-number-rate-distortion} 中下界证明的关键成分是信息论中的一个重要结果，称为失真编码率的\textit{香农下界}（Shannon lower bound），以克劳德·香农（Claude Shannon）的名字命名，他首先在一个特例中推导出了它 \cite{shannon1959coding}。它断言了对于任何具有密度 $p(\vx)$ 和有限期望平方范数的随机变量 $\vx$，失真编码率函数的以下估计 \cite{Linder1994-ej}：
    \begin{equation}\label{eq:slb}
        \cR_{\epsilon}(\vx)
        \geq
        h(\vx)
        - \log_2 \volume(B_{\epsilon})
        -
        C_D,
    \end{equation}
    其中 $C_D > 0$ 是仅依赖于 $D$ 的常数。此外，当 $\epsilon \to 0$ 时，这个下界实际上是紧的：即，
    \begin{equation}
        \lim_{\epsilon \to 0} \cR_{\epsilon}(\vx) - \left[ h(\vx) - \log_2
        \volume(B_{\epsilon}) - C_D\right] = 0.
    \end{equation}
    因此，当失真 $\epsilon$ 很小时，我们可以主要从香农下界的角度来思考，而不必直接处理定义失真编码率的（通常难以求解的）优化问题 \eqref{eqn:rate-distortion-general}。

    香农下界是连接编码率、熵最小化/去噪以及用于学习低维分布的几何球堆积方法的桥梁。请注意，在均匀密度 $p(\vx)$ 的特殊情况下，\eqref{eq:slb} 变为
    \begin{align}
        \cR_{\epsilon}(\vx) &\geq -\int_K \frac{1}{\volume(K)} \log_2
        \frac{1}{\volume(K)} \odif \vxi
        - \log_2 \volume(B_{\epsilon}) - C_D
        \\
        &=
        \log_2 \volume(K) / \volume(B_{\epsilon})
        - C_D. \label{eq:shannon-lower-bound-uniform}
    \end{align}
    比率 $\volume(K) / \volume(B_{\epsilon})$ 通过最坏情况的论证近似了覆盖 $K$ 所需的 $\epsilon$-球数量，当 $\epsilon$ 很小时，这对于足够规则的集合 $K$ 是准确的（详见\Cref{app:rate-distortion-covering}）。同时，回想一下本章前面的高斯去噪模型 $\vx_{\epsilon} = \vx + \epsilon \vg$，其中 $\vg \sim \cN(\vzero, \vI)$ 独立于 $\vx$。有趣的是，联合分布 $(\vx, \epsilon\vg)$ 的微分熵可以计算为
    \begin{align}
        h(\vx, \epsilon\vg)
        &=
        -\int
        p(\vxi) p(\vgamma) \log_2 p(\vxi) p(\vgamma) \odif \vxi \odif \vgamma
        \\
        &=
        h(\vx) + h(\epsilon \vg).
    \end{align}
    我们已经看到了在\Cref{eqn:entropy-Gaussian-multi}中计算的高斯熵：当 $\epsilon$ 很小时，它等于（在相差一个加性常数的意义下）我们在香农下界中看到的体积量 $-\log_2 \volume(B_{\epsilon})$。在某些特殊情况下（例如，支撑在非相干低秩子空间上的数据），当 $\epsilon$ 很小且 $\vx$ 的支撑足够规则时，$\vx_{\epsilon}$ 的分布甚至可以在局部被分布 $p(\vx)$ 和 $p(\vg)$ 的乘积很好地近似，从而证明了上述计算的合理性。因此，高斯去噪过程产生了对香农下界的另一种解释，即它源于 $\vx$ 的噪声版本的熵，其噪声水平与失真水平 $\epsilon$ 成正比。

    因此，这种基于球覆盖的有限失真编码率方法统一并推广了以往对分布复杂度的度量，使我们能够以统一方式区分并排序不同分布。这些相互关联的观点在\Cref{fig:continuation}中得到了可视化展示。
\end{remark}

对于有限失真水平下的一般分布，通常难以得到失真编码率函数的解析形式。人们往往必须求助于数值计算\footnote{感兴趣的读者可以参考 \cite{Blahut-1972}，了解一种对离散分布的失真编码率函数进行数值计算的经典算法。}。尽管如此，正如我们将看到的，在我们的语境下，我们通常需要把失真编码率写成数据点集合或其表征的显式函数。这是因为我们希望以编码率作为衡量表征优劣的指标。显式解析形式使我们更容易判断如何变换数据分布以改进表征。因此，我们将聚焦于那些失真编码率函数具有显式解析形式的分布。为此，我们从最简单且最重要的分布族开始。

\subsection{低维高斯分布的有损编码率}\label{subsec:lossy DR}
现在假设我们给定了一组来自任意分布的数据样本 $\vX = [\vx_1, \ldots, \vx_N]$。\footnote{或者这些数据点本身可以被视为一种（经验）分布。}我们希望提出一种建设性的方案，能够以一定的精度对数据进行编码，例如
\begin{equation}
    \vx_i \mapsto \hat \vx_i, \quad \text{s.t.} \quad \|\vx_i -
    \hat \vx_i\|_2 \le \epsilon.
\end{equation}
请注意，这是一个充分、明确且可解释的条件，它确保数据编码后满足 \(\frac{1}{N}\sum_{i = 1}^{N} \norm{\vx_i- \hat \vx_i}_2^{2} \le \epsilon^{2}\)。后者不等式正是给定经验分布及其编码所对应的失真编码率约束。例如，在\Cref{example:sphere-covering-rate-distortion}中，我们使用这一简化标准来显式求解给定编码率下的最小失真和对应编码方案。

不失一般性，让我们假设 $\vX$ 的均值为零，即 $\frac{1}{N} \sum_{i = 1}^{N} \vx_i = \boldsymbol{0}$。在没有任何关于 $\vX$ 背后分布性质的先验知识的情况下，我们可以将 $\vX$ 视为从具有以下协方差的高斯分布 $\mathcal{N}(\boldsymbol{0}, {\boldsymbol{\Sigma}})$ 中采样的\footnote{众所周知，在给定固定方差的情况下，高斯分布达到最大熵。也就是说，它给出了可能编码率的最坏情况的上限。}
\begin{equation}
    {\boldsymbol{\Sigma}} = \frac{1}{N} \vX\vX^\top.
\end{equation}
请注意，在几何上 ${\boldsymbol{\Sigma}}$ 表征了大多数样本 $\vx_i$ 所在的椭球区域。

我们可以将 $\hat \vX = [\hat \vx_1,\ldots, \hat \vx_N]$ 视为 $\vX = [\vx_1, \ldots, \vx_N]$ 的噪声版本：
\begin{equation}
    \hat \vx_i = \vx_i + \vw_i,
\end{equation}
其中 $\vw_i  \sim
\mathcal{N}(\boldsymbol{0} , {\epsilon^2}  \boldsymbol{I}/{D})$ 是独立于 $\vx_i$ 的高斯噪声。那么 $\hat \vx_i$ 的协方差由下式给出
\begin{equation}
    \hat{\boldsymbol{\Sigma}} = \mathbb{E}\left[\hat \vx_i \hat
    \vx_i^\top\right] = \frac{\epsilon^2}{D} \boldsymbol{I} +
    \frac{1}{N} \vX\vX^\top.
\end{equation}
请注意，由向量 $\vx_i$ 张成的区域的体积与协方差矩阵行列式的平方根成正比
\begin{equation}
    \volume(\hat \vx_i) \propto \sqrt{\det
    \big(\hat{\boldsymbol{\Sigma}}\big)} =
    \sqrt{\det\left(\frac{\epsilon^2}{D} \boldsymbol{I} + \frac{1}{N}
    \vX\vX^\top \right)}.
\end{equation}
每个随机向量 $\boldsymbol{w}_i$ 张成的体积与以下成正比
\begin{equation}
    \volume(\boldsymbol{w}_i) \propto
    \sqrt{\det\left(\frac{\epsilon^2}{D} \boldsymbol{I} \right)}.
\end{equation}

\begin{figure}
    \centering
    \includegraphics[width=\linewidth]{\toplevelprefix/chapters/chapter4/figs/Gaussian-compression.png}
    \caption{使用 $\epsilon$-球覆盖数据向量张成的区域。空间的体积越大，需要的球就越多，因此需要更多的比特来编码和枚举这些球。每个实值向量 $\vx$ 都可以编码为它落入的球的编号。}
    \label{fig:ball-packing}
\end{figure}

为了对落入 $\hat \vx_i$ 张成区域的向量进行编码，我们可以用半径为 $\epsilon$ 的非重叠球覆盖该区域，如\Cref{fig:ball-packing}所示。当 $\hat \vx_i$ 张成区域的体积明显大于 $\epsilon$-球的体积时，我们覆盖该区域所需的球的总数大约等于这两个体积的比值：
\begin{equation}
    \#\,\epsilon\text{-球} \approx \frac{\volume(\hat
    \vx_i)}{\volume(\vw_i)} = \sqrt{\det\left(\boldsymbol{I} +
    \frac{D}{N\epsilon^2} \vX\vX^\top  \right)}.
\end{equation}
如果我们使用二进制数来标记感兴趣区域中的所有 $\epsilon$-球，那么所需的二进制比特总数因此为 %
\begin{equation}
    \cR_{\epsilon}(\vX) \approx %
    \log_2 (\#\,\epsilon\text{-球}) \approx R_{\epsilon}(\vX)
    \doteq \frac{1}{2} \log \det \left(\boldsymbol{I} +
    \frac{D}{N\epsilon^2} \vX\vX^\top \right).
    \label{eqn:rate-Gaussian}
\end{equation}
请注意，由\Cref{ex:gaussian_rate_distortion}可知，该公式右侧与协方差矩阵为 \(\frac{1}{N} \vX\vX^\top + (\epsilon^{2}/D) \vI\)（即数据经验协方差加噪声项）的高斯信源在精度 \(\epsilon\) 下的失真编码率一致。这样，关于样本张成体积与球堆积的几何论证就转化为关于受扰动高斯信源失真编码率的信息论论证。

\begin{example}
    \Cref{fig:ball-packing} 显示了一个具有椭圆支撑的二维分布示例——可近似看作二维高斯分布的支撑。该区域被半径为 $\epsilon$ 的小球覆盖。设这些球的编号为 $1$ 到 $n$。给定该区域中的任意向量 $\vx$，我们只需确定它最接近哪个 $\epsilon$-球中心，记为 $\operatorname{ball}_{\epsilon}(\vx)$。要存储 $\vx$，只需存储该球的编号，这需要 $\log(n)$ 个比特。若要从该编号解码 $\vx$，我们只需将 $\hat \vx$ 取为该球中心。这给出了一个显式的编码-解码方案：
    \begin{equation}
        \vx \longrightarrow \operatorname{ball}_{\epsilon}(\vx)
        \longrightarrow \hat \vx = \text{中心}\big(\operatorname{ball}_{\epsilon}(\vx)\big).
    \end{equation}
    这些球心可视为编码方案中码本（或字典）的“码字”。容易看出，这种（有损）编码-解码方案的精度约为球半径 $\epsilon$。显然，$\cR_{\epsilon}(\vZ)$ 是使用该编码方案对每个向量 $\vz$ 的球编号进行编码所需的平均比特数，因此可称为该方案对应的{\em 编码率}。
\end{example}

从上面的推导中，我们知道编码率 $\cR_{\epsilon}(\vX)$ 是可以通过显式编码（和解码）方案（近似）实现的。它有两个有趣的性质：
\begin{itemize}
    \item 首先，人们可能会注意到 $R_{\epsilon}(\vX)$ 非常类似于高斯信源的失真编码率函数 \cite{Cover-Thomas}（见\Cref{ex:gaussian_rate_distortion}）。事实上，当 $\epsilon$ 很小时，上述表达式是高斯信源失真编码率的近似，正如 \cite{MaY2007-PAMI} 所指出的那样。
    \item 其次，如果假设数据 $\vX$ 来自线性子空间，则可以推导出相同的闭式编码率 $R_{\epsilon}(\vX)$ 作为 \(\cR_{\epsilon}(\vX)\) 的近似。这可以通过适当地量化 $\vX = \boldsymbol{U} \boldsymbol{\Sigma}\boldsymbol{V}^\top$ 的奇异值分解（SVD）并为 $\boldsymbol{U}$ 张成的子空间中的向量构建有损编码方案来证明 \cite{MaY2007-PAMI}。
\end{itemize}
在我们的背景下，闭式表达式 $R_{\epsilon}(\vX)$ 是相当基础的：它是与从高斯分布或线性子空间中提取的数据的显式且自然的有损编码方案相关的编码率。正如我们将在下一章中看到的，这个公式在理解深度神经网络的架构中起着重要作用。

\subsection{低维高斯混合的聚类}
\label{sec:clustering-Gaussians}
正如我们之前讨论的，给定数据集 $\vX$ 往往具有低维内在结构。因此，将其编码为一般高斯分布会非常冗余。若能识别出 $\vX$ 的这些内在结构，就可以设计更好的编码方案，从而显著降低编码率。等价地，用于编码 $\vX$ 的码可以进一步压缩。我们将看到，压缩为识别此类结构提供了一种统一且可计算的方法。本节将以最基本的低维结构族来说明这一思想：低维高斯分布或子空间的混合。

\begin{example}
    \Cref{fig:two-subspaces} 显示了一个示例，其中数据 $\vX$ 分布在两个子空间（或低维高斯分布）周围。如果将它们视为一个单一的高斯分布并一起编码，由所有蓝球表示的关联离散（有损）码本显然是非常冗余的。我们可以尝试识别这两个子空间的位置（记为 $S_1$ 和 $S_2$），并设计一个仅覆盖这两个子空间的码本，即绿球。如果我们能正确地将数据 $\vX$ 中的样本划分到这两个子空间中：$\vX = [\vX_1, \vX_2]\bm \Pi$，其中 $\vX_1 \in S_1$ 且 $\vX_2 \in S_2$，其中 $\bm \Pi$ 表示置换矩阵，那么由此产生的数据编码率将低得多。这提供了一种更简约、因此更理想的数据表征。
\end{example}

\begin{figure}
    \centering
    \includegraphics[width=0.5\linewidth]{\toplevelprefix/chapters/chapter4/figs/Two-subspaces.png}
    \caption{分布在两个子空间周围的数据的两种有损编码方案的比较。一种是为两个子空间张成的整个空间堆积（蓝色）$\epsilon$-球；另一种是仅在两个子空间周围的管状邻域内堆积球。后者显然具有小得多的码本，并导致子空间上样本的编码率低得多。}
    \label{fig:two-subspaces}
\end{figure}

因此，更一般地说，如果数据是从子空间或低维高斯分布的任何混合中提取的，那么识别这些分量并根据这些分量的内在维度对数据进行编码将是理想的。事实证明，假设数据是从低维高斯混合中提取的，我们并没有失去太多的普遍性。这是因为高斯混合可以紧密近似大多数一般分布 \cite{borkar2016gaussian}。

\paragraph{聚类问题。}
现在对于这个特定的分布族，我们如何从一组样本中有效且高效地识别出这些低维分量
\begin{equation}
    \vX = \left[\vx_1, \vx_2, \ldots, \vx_N\right],
\end{equation}
？换句话说，给定整个数据集 $\vX$，我们希望将其划分或聚类为多个（例如 $K$ 个）子集：
\begin{equation}
    \vX\bm \Pi = [\vX_1, \vX_2, \dots, \vX_K],
\end{equation}
其中每个子集由仅从一个低维高斯分布或子空间中提取的样本组成，$\bm \Pi$ 是一个置换矩阵，用于指示划分的成员归属。请注意，根据情况，划分可以是确定性的，也可以是概率性的。正如 \cite{ma2007segmentation} 所示，对于高斯混合，概率性划分不会导致更低的编码率。因此，为简单起见，我们这里仅考虑确定性划分。

\paragraph{通过有损压缩进行聚类。}
解决上述聚类问题的主要困难在于，我们通常不知道簇的数量 $K$，也不知道每个分量的维度。对这个聚类问题的研究有着悠久的历史。教科书 \cite{GPCA} 对解决该问题的不同方法进行了系统而全面的介绍。为了找到解决这个问题的有效方法，我们首先需要理解并澄清我们为什么要聚类。换句话说，与不聚类相比，我们从数据聚类中到底能得到什么？我们如何衡量这种收益？从数据压缩的角度来看，正确的聚类应该导致更高效的编码（和解码）方案。

对于任何给定的数据集 $\vX$，已经有两个明显的编码方案作为基线。它们代表了两种极端的编码数据的方式：
\begin{itemize}
    \item 简单地将所有样本视为从一个单一的高斯分布中提取的。如前所述，相关的编码率由下式给出：
        \begin{equation}
            \cR_{\epsilon}(\vX) \approx R_{\epsilon}(\vX) =
            \frac{1}{2} \log \det \left(\boldsymbol{I} +
            \frac{D}{N\epsilon^2} \vX\vX^\top \right).
        \end{equation}
    \item 简单地通过为每个样本分配不同的编号来分别记住所有样本。编码率将是：
        \begin{equation}
            \cR_0(\vX) = \log(N).
        \end{equation}
\end{itemize}

请注意，对于量化误差 $\epsilon$ 的某些（极端）选择，这两种编码方案都可以成为“最优”解：
\begin{enumerate}
    \item {\em 懒惰机制}：如果我们选择极大的 $\epsilon$，$\vX$ 中的所有样本都可以被一个球覆盖。速率为  $\lim_{\epsilon \rightarrow \infty} \cR_\epsilon
        \rightarrow \frac{1}{2}\log\det (\boldsymbol{I}) = 0$。
    \item {\em 记忆机制}：如果 $\epsilon$ 极小，$\vX$ 中的每个样本都被不同的 $\epsilon$-球覆盖，因此总数为 $N$。速率为
        $\lim_{\epsilon \rightarrow 0} \cR_\epsilon \rightarrow \log(N)$。
\end{enumerate}
请注意，第一种方案对应于人们根本不关心分布的任何有趣特征的场景。人们不想为任何有信息量的内容腾出任何比特。我们称之为“懒惰机制”。第二种方案对应于人们想要以极高精度解码每个样本的场景。所以人们最好“记住”每一个样本。我们称之为“记忆机制”。

请注意，这些关于 \(\epsilon\) 的渐近性正是我们在\Cref{sec:diffusion_memorization}末尾引入有损编码率作为实际训练的扩散模型所近似的量时所讨论的。在当前情况下，我们\textit{显式地}将分布编码到精度 \(\epsilon\)；在扩散模型的情况下，我们\textit{隐式地}编码分布。该方案的某些属性（如记忆和泛化）在两种情况下表现相似，实际上人们可以通过研究一种情况来理解另一种情况。

\begin{figure}
    \centering
    \includegraphics[width=0.9\linewidth]{\toplevelprefix/chapters/chapter4/figs/circle-packing.png}
    \caption{二维平面上的若干随机样本。考虑为每个样本分配一个以该样本为中心的 $\epsilon$-圆盘。样本的密度从左到右增加。}
    \label{fig:circle-packing}
\end{figure}
\begin{example}
    为了了解何时倾向于使用记忆机制，让我们考虑在二维平面的单位面积内随机分布的若干（例如 $N$ 个）样本。\footnote{假设这些点是由密度为每单位面积 $N$ 个点的泊松过程提取的。}想象一下，我们试图设计一个具有固定量化误差 $\epsilon$ 的有损编码方案。这相当于在每个样本周围放置一个 $\epsilon$-圆盘，如\Cref{fig:circle-packing}所示。当 $N$ 很小时，所有圆盘相互重叠的几率为零。在这种情况下，大小为 $N$ 的码本是必要且最优的。当 $N$ 或密度达到某个临界值 $N_c$ 时，所有圆盘很有可能开始重叠并连接成一个覆盖整个平面的簇——这种现象被称为连续“渗流”（percolation） \cite{Gilbert-1961,Mertens-Moore-2012}。当 $N$ 变得大于这个值时，圆盘会严重重叠。圆盘的数量 $N$ 变得非常冗余，因为我们只想在给定精度 $\epsilon$ 内对平面上的点进行编码。覆盖所有样本所需的圆盘数量远小于 $N$。\footnote{事实上，有高效的算法可以找到这样的覆盖 \cite{Booth-2001}。}
\end{example}

懒惰机制和记忆机制都在某种程度上是微不足道的，也许没有什么理论或实际意义。当用于编码从具有{\em 紧致且低维支撑}的分布中提取的大量样本时，这两种方案都远非最优。有趣的机制存在于这两者之间。
\begin{example}
    \begin{figure}[t]
        \centering
        \includegraphics[width=0.4\linewidth]{\toplevelprefix/chapters/chapter4/figs/Two-lines-and-plane.png}
        \includegraphics[width=0.9\linewidth]{\toplevelprefix/chapters/chapter4/figs/Coding-Rate.jpg}
        \caption{上：从 $\R^3$ 中的两条直线和一个平面中提取的 358 个噪声样本。下：改变 $\epsilon$ 对聚类结果和编码率的影响。红线标记了添加到样本中的高斯噪声的方差 $\epsilon_0$。}
        \label{fig:two-lines-and-plane}
        \label{fig:two-lines-and-plane-epsilon}
    \end{figure}
    \Cref{fig:two-lines-and-plane} 给出了一个示例：噪声样本来自 $\R^3$ 中的两条直线和一个平面。正如右侧图 (c) 所示，随着 $\epsilon$ 增大，最优编码率单调下降，这与失真编码率函数的性质一致。图 (a) 和 (b) 展示了当 $\epsilon$ 从极小（接近零）变化到极大（趋向无穷）时，使编码率最小的最优簇数。我们可以清楚地看到图两端对应懒惰机制和记忆机制。同时，从图 (b) 也可以看出：当量化误差 $\epsilon$ 选在真实噪声方差 $\epsilon_0$ 附近时，最优簇数为“正确”的 3，对应两条直线和一个平面。我们将这一中间机制非正式地称为“泛化机制”。请注意，这些机制之间会发生急剧的相变。\footnote{据我们所知，这些相变行为目前尚无严格理论证明。}
\end{example}

从上述讨论和例子中，我们看到：当量化误差与样本密度\footnote{或样本密度与量化误差}处于适当范围时，最小化有损编码率能够揭示采样数据背后的潜在（低维）分布。{\em 因此，量化最初虽然是出于实用性的选择，但现在似乎正在成为从有限样本的经验分布中学习连续分布的必要条件。}尽管解释这一现象的严格理论仍不完善，但在这里，出于学习目的，我们更关心如何利用这一现象来设计能够恢复正确分布的算法。

让我们使用\Cref{fig:two-subspaces}中所示的简单示例来说明基本思想。如果能将 $\vX$ 中的所有样本划分为 $\vX_1$ 和 $\vX_2$ 两个簇，分别有 $N_1$ 和 $N_2$ 个样本，那么相关的编码率将是\footnote{我们在这里忽略了编码每个样本的成员归属所需的一些开销比特，例如通过霍夫曼编码。}
\begin{equation}
    R_{\epsilon}^c(\vX\mid \boldsymbol{\Pi}) =
    \frac{N_1}{N}R_{\epsilon}(\vX_1) + \frac{N_2}{N}R_{\epsilon}(\vX_2),
\end{equation}
其中我们使用 $\boldsymbol{\Pi}$ 来指示划分的成员归属。如果划分尊重分布的低维结构，在这种情况下 $\vX_1$ 和 $\vX_2$ 分别属于两个子空间，那么由此产生的编码率应该显著小于上述两个基本方案：
\begin{equation}
    R_{\epsilon}^c(\vX \mid \boldsymbol{\Pi}) \ll R_{\epsilon}(\vX),
    \quad     R_{\epsilon}^c(\vX \mid \boldsymbol{\Pi}) \ll R_0(\vX).
\end{equation}
一般来说，我们可以将聚类问题转化为最小化编码率的优化问题：
\begin{equation}
    \min_{\boldsymbol{\Pi}}  \bc{ R_{\epsilon}^c(\vX \mid \boldsymbol{\Pi})
    \doteq \sum_{k=1}^K \frac{N_k}{N}R_{\epsilon}(\vX_k)}.
\end{equation}

\paragraph{聚类的优化策略。}
剩下的问题是如何优化上述编码率目标以找到最优簇。实现这一目标有三种自然的方法：
\begin{enumerate}
    \item 从将整个集合 $\vX$ 视为单一聚类开始（即懒惰机制），再搜索（例如随机搜索）其划分，使编码率变小。
    \item 反过来，从将每个样本 $\vx_i$ 视为独立聚类开始（即记忆机制），再搜索可使编码率下降的聚类合并。
    \item 或者，如果能将隶属度 $\boldsymbol{\Pi}$ 表示（或近似）为连续参数，就可使用梯度下降（GD）等优化方法。
\end{enumerate}
第一种方法在计算上并不吸引人，因为需要尝试的划分数量随样本数呈指数增长。例如，将 $\vX$ 划分为两个大小相等子集的划分数为 $N \choose N/2$，当 $N$ 增大时会迅速爆炸。我们将在下一章 \Cref{ch:deep-networks} 中探讨第三种方法。在那里我们将看到，深度神经网络（尤其是 Transformer）的作用如何与编码率目标建立联系。

第二种方法最初由 \cite{ma2007segmentation} 提出。它展示了高效评估编码率（例如使用解析形式）的价值。借助这一点，可以通过最小化编码长度（MCL）原理，高效且有效地找到数据的（低维）聚类。请注意，对于包含 $N_k$ 个样本的聚类 $\vX_k$，编码 $\vX_k$ 中所有样本所需的二进制比特长度为：\footnote{事实上，对编码长度更准确的估计是 $L(\vX_k) = (N_k+D) R_\epsilon(\vX_k)$，其中额外比特用于编码子空间的基 \cite{ma2007segmentation}。为简洁起见，我们在此省略该开销。}
\begin{equation}
    L(\vX_k) = N_k R_\epsilon(\vX_k).
\end{equation}
如果给定两个聚类 $\vX_k$ 和 $\vX_l$，并希望将样本编码为两个独立聚类，则所需的二进制比特长度为
\begin{equation*}
    L^c(\vX_k, \vX_l) = N_k R_\epsilon(\vX_k) + N_l R_\epsilon(\vX_l) -
    N_k \log\frac{N_k}{N_k + N_l} - N_l \log\frac{N_l}{N_k + N_l}.
\end{equation*}
最后两项是根据哈夫曼编码（Huffman code）对样本的隶属度进行编码所需的比特数。

然后，给定任意两个独立聚类 $\vX_k$ 和 $\vX_l$，我们可以根据两个编码长度之差来决定是否合并它们：
\begin{equation}
    L(\vX_k \cup \vX_l) - L^c(\vX_k, \vX_l)
\end{equation}
其符号决定是否合并，其中 $\vX_k \cup \vX_l$ 表示 $\vX_k$ 和 $\vX_l$ 的样本并集。若该值为负，意味着将这两个聚类合并后编码长度会变小。这个简单事实引出了 \cite{ma2007segmentation} 提出的如下聚类算法：
\begin{algorithm}[!htbp]
    \caption{编码长度的成对最速下降法}\label{alg:steepest_descent_coding_length}
    \begin{algorithmic}[1]
        \Require{\(N\) 个数据点 \(\{\vx_{i}\}_{i = 1}^{N}\)}
        \Ensure{聚类集合 \(\cC\)}

        \Procedure{PairwiseSteepestDescentOfCodingLength}{$\{\vx_{i}\}_{i
        = 1}^{N}$}
        \State{\(\cC \gets \{\{\bm x_{i}\}\}_{i = 1}^{N}\)}
        \Comment{初始化 \(N\) 个各包含一个元素的聚类 \(\vX_{k}\)}
        \While{\(\abs{\cC} > 1\)}
        \If{\(\displaystyle \min_{\vX_{k}, \vX_{l} \in \cC}[L(\vX_{k}
        \cup \vX_{l}) -L^{c}(\vX_{k}, \vX_{l})] \geq 0\)} \Comment{如果任何合并都无法节省比特}
        \State{\Return{\(\cC\)}} \Comment{提前返回 \(\cC\) 并退出}
        \Else
        \State{\(\displaystyle \vX_{k^{\ast}}, \vX_{l^{\ast}} \gets
                \argmin_{\vX_{k}, \vX_{l} \in \cC}[L(\vX_{k} \cup \vX_{l}) -
        L^{c}(\vX_{k}, \vX_{l})]\)} \Comment{合并能节省最多比特的聚类}
        \State{\(\displaystyle \cC \gets [\cC \setminus
                \{\vX_{k^{\ast}}, \vX_{l^{\ast}}\}] \cup \{\vX_{k^{\ast}}
        \cup \vX_{l^{\ast}}\}\)} \Comment{移除未合并的聚类并加回合并后的聚类}
        \EndIf
        \EndWhile
        \State{\Return{\(\cC\)}}  \Comment{如果所有合并都能带来节省，则返回一个聚类}
        \EndProcedure
    \end{algorithmic}
\end{algorithm}

请注意，该算法是可处理的，因为（成对）比较与合并的总次数约为 $O(N^2\log N)$。然而，由于其贪婪性质，尚无理论保证该过程会收敛到全局最优聚类解。尽管如此，正如 \cite{ma2007segmentation} 所报告的，这个看似简单的算法在实践中效果极好。\Cref{fig:two-lines-and-plane} 中的聚类结果实际上就是由该算法计算得到的。

\begin{example}[图像分割]\label{eg:image-segmentation} 上述编码长度度量及其相关聚类算法假设数据分布是（低维）高斯混合。尽管这一假设看起来略显理想化，但在模型近似成立的场景中，该度量和算法已经非常有用，甚至很强大。

    例如，自然图像通常由多个具有几乎均匀纹理的区域组成。如果我们从每个区域中提取许多小窗口，它们应该类似于从（低维）高斯分布中提取的样本，如 \Cref{fig:image-patch} 所示。\Cref{fig:image-segmentation} 展示了将上述聚类算法直接应用于图像块的图像分割结果。关于针对图像分割问题定制该算法的更多技术细节，可以在 \cite{Mobahi-IJCV2011} 中找到。
\end{example}

\begin{figure}
    \centering
    \includegraphics[width=0.4\linewidth]{\toplevelprefix/chapters/chapter4/figs/image-segmentation-tiles.png}
    \caption{大小为 $w\times w$ 像素的图像块。}
    \label{fig:image-patch}
\end{figure}

\begin{figure}[th]
    \centering
    \includegraphics[width=0.8\linewidth]{\toplevelprefix/chapters/chapter4/figs/image-segmentation.png}
    \caption{基于应用于图像块的聚类算法的分割结果。}
    \label{fig:image-segmentation}
\end{figure}

\section{通过最大化信息增益进行变换}
\label{sec:chap4-representation-learning-problem}

到目前为止，本章已经讨论了如何通过压缩原理识别具有低维结构的分布。正如前两节所示，计算意义上的压缩可以通过去噪或聚类来实现。\Cref{fig:Gaussian-Subspaces} 用一个典型示例说明了这一点。

当然，识别数据分布的最终目的是服务于后续任务，例如（图像）分割、分类或生成。因此，所得分布如何被“表征”，对于高效检索和利用与下游任务相关的信息至关重要。这自然引出一个基本问题：{\em 什么样的表征对下游任务才是真正“好”的？} 下文将探讨有意义且有用的表征应具备的基本属性，以及如何通过最大化信息增益来刻画并追求这些属性。

\paragraph{如何衡量表征的好坏。}
可以将给定数据集视为高维空间（例如 $\R^D$）中随机向量 $\vx$ 的样本，其分布通常具有远低于环境空间维度的内在维度。一般而言，{\em 学习表征}（learning a representation）是指学习一个连续映射（如 $f(\cdot)$），将 $\vx$ 变换到另一个（通常更低维的）空间（如 $\R^d$，其中 $d < D$）中的 {\em 特征向量}（feature vector）$\vz$。我们希望通过如下映射
\begin{equation}
    \vx \in \R^D \xrightarrow{\hspace{2mm} f(\vx)\hspace{2mm}}
    \vz  \in \R^d,
    \label{eqn:chap4-1-encoding}
\end{equation}
$\vx$ 的低维内在结构能够被识别，并由 $\vz$ 以更紧凑、更结构化的方式表征，从而促进分类、生成等后续任务。特征 $\vz$ 可视为原始数据 $\vx$ 的（学习到的）紧凑编码，因此映射 $f$ 也被称为 \textit{编码器}（encoder）。表征学习的基本问题是
\begin{center}
    \noindent{\em 衡量表征好坏的原则性且有效的标准是什么？}
\end{center}

\begin{figure}[t]
    \centering
    \includegraphics[width=0.8\linewidth]{\toplevelprefix/chapters/chapter4/figs/Gaussian-Subspaces.png}
    \caption{从一般的随机高斯分布（右）开始，通过去噪或聚类识别具有两个子空间的低维分布（左）。}
    \label{fig:Gaussian-Subspaces}
\end{figure}

从概念上看，表征 $\vz$ 的质量取决于两点：它在多大程度上保留了 $\vx$ 中与后续任务最相关、最充分的信息，以及它表征这些信息的效率。长期以来，人们常将学习到的表征的“充分性”或“好坏”定义为特定任务相关。例如在分类中，$\vz$ 只需足以预测类别标签 $\y$。下面我们从经典图像分类问题出发，说明这种任务特定的“表征”概念存在局限，因而需要推广。

\subsection{从分类数据中学习表征}

假设 $\bm{x} \in \R^D$ 是从 $K$ 个分量分布的混合 $\mathcal{D} = \{\mathcal{D}_k\}_{k=1}^K$ 中抽取的随机向量。给定随机向量 $\bm x$ 的有限独立同分布（i.i.d.）样本集 $\vX = [\vx_1, \vx_2, \ldots, \vx_N] \in \Re^{D\times N}$，我们希望通过连续映射 $f(\vx): \R^D \rightarrow \R^d$ {\em 学到一个好的表征}，使其既能捕捉 $\vx$ 的内在结构，又能尽可能促进后续分类任务。\footnote{分类是深度学习最初取得成功的领域，引发了人们对深度网络的爆炸性兴趣。尽管我们的研究侧重于分类，但我们相信这些思想和原则可以自然推广到其他设置，例如回归。} 为了简化分布 $\mathcal{D}$ 的学习，在常见的监督分类设置中，每个样本 $\vx_i$ 都带有真实类别标签（或类别码字），通常表示为独热（one-hot）向量 $\y_i \in \R^K$。

\paragraph{通过交叉熵编码类别信息。} 大量研究表明，对于许多实际数据集（例如图像、音频和自然语言），从数据 $\bm{x}$ 到类别标签 $\bm{y}$ 的（编码）映射可以通过训练深度网络有效建模，\footnote{这里暂不讨论具体应使用哪类网络及其原因。此处只需考虑任意经实证验证有效的深度网络；网络架构的合理性将留到下一章讨论。} 记作 $$f(\vx, \theta):\vx \mapsto \y$$，其中网络参数 $\theta \in \Theta$，$\Theta$ 为参数空间。为使输出 $f(\vx, \theta)$ 与标签 $\y$ 良好匹配，我们希望在训练集 $\{(\vx_i, \y_i)\}_{i=1}^N$ 上最小化 {\em 交叉熵损失}（cross-entropy loss）：
\begin{equation}
    \min_{\theta \in \Theta} \; - \mathbb{E}[\langle \y, \log(f(\vx,
    \theta)) \rangle] \, \approx - \frac{1}{N}\sum_{i=1}^N \langle
    \y_i, \log\left(f(\vx_i, \theta)\right) \rangle.
    \label{chap4-eqn:cross-entropy}
\end{equation}
最优网络参数 $\theta$ 通常通过高效的梯度下降方案优化上述目标来找到，其中梯度通过反向传播计算，如 \Cref{app:optimization} 的 \Cref{app:BP-section} 中所述。

尽管该方法有效且广受欢迎，但仍有两个严重局限：（1）它只追求预测标签 $\y$，即使标签可能被错误标注。实证研究表明，作为“黑盒”使用的深度网络甚至可以拟合随机标签~\cite{zhang2017understanding}。（2）在这种端到端数据拟合框架下，尽管已有大量实证工作尝试解释学习到的特征，%
我们仍不清楚网络中间特征在多大程度上捕捉了使分类得以成立的数据内在结构。学习特征的精确几何与统计属性也往往被掩盖，进而导致深度学习缺乏可解释性和后续性能保证（如泛化性、可迁移性、鲁棒性等）。因此，{\em 本节目标之一是将目标重述为学习数据 $\vx$ 的明确、有意义且有用的表征（不限于分类），以克服这些局限。}

\begin{figure}
    \centering
    \includegraphics[width=0.95\linewidth]{\toplevelprefix/chapters/chapter4/figs/neural_collapse.pdf}
    \caption{在包含 3 个随机选择类别的 CIFAR10 数据集上训练时，VGG13 神经网络倒数第二层输出的演变。图片来自 \cite{papyan2020prevalence}。}
    \label{chap4-fig:neural-collapse}
\end{figure}

\paragraph{通过信息瓶颈获取最小判别特征。}
解释深度网络作用的一种流行观点是：将网络中间层输出视为从数据中提取的潜在特征 $\vz = f(\vx, \theta) \in \Re^d$，这些特征在多个类别之间具有判别性。随后，通过优化分类器 $g(\vz)$，学习到的表征 $\vz$ 可用于预测类别标签 $\y$：
\begin{equation}
    \vx   \xrightarrow{\hspace{2mm} f(\vx, \theta)\hspace{2mm}} \vz
    \xrightarrow{\hspace{2mm} g(\vz) \hspace{2mm}} \y.
\end{equation}
我们从 \eqref{eq:mutual} 中知道，两个随机变量（例如 $\vx,\vz$）之间的 {\em 互信息}（mutual information）定义为
\begin{equation}
    I(\vx; \vz) = H(\vx) - H(\vx\mid \vz),
\end{equation}
其中 $H(\vx \vert \vz)$ 是给定 $\vz$ 时 $\vx$ 的条件熵。互信息也称为 {\em 信息增益}（information gain）：它衡量给定 $\vz$ 后，随机变量 $\vx$ 的熵可减少多少；等价地，它衡量 $\vz$ 包含了多少关于 $\vx$ 的信息。{\em 信息瓶颈}（information bottleneck, IB）框架 \cite{Tishby-ITW2015} 进一步假设网络的作用是学习 $\vz$，使其成为预测 $\y$ 的最小充分统计量。形式上，它寻求最大化 $\vz$ 与 $\y$ 之间的互信息 $I(\vz, \y)$，同时最小化 $\vx$ 与 $\vz$ 之间的互信息 $I(\vx, \vz)$：
\begin{equation}
    \max_{\theta\in \Theta}\; \mbox{IB}(\vx, \y, \vz) \doteq I(\vz; \y) -
    \beta I(\vx; \vz) \quad\ \mathrm{s.t.}\ \vz = f(\vx, \theta),
    \label{chap4-eqn:information-bottleneck}
\end{equation}
其中 $\beta >0$。

如果能够克服该框架的一些技术注意事项 \cite{kolchinsky2018caveats-ICLR2018}（例如如何用来自退化分布的有限样本准确评估互信息），它确实有助于解释深度网络的某些行为。比如，近期工作 \cite{papyan2020prevalence} 表明，通过交叉熵损失 \eqref{chap4-eqn:cross-entropy} 学到的表征会呈现 \emph{神经坍缩}（neural collapse）现象：每个类别的特征被映射到一个一维向量，而该类别中的其他信息被抑制，如 \Cref{chap4-fig:neural-collapse} 所示。\footnote{我们稍后将在第 \ref{eg:Rate-Reduction-CIFAR10} 节的真实图像实验中观察到这一现象。}
\begin{remark}[神经坍缩]
    神经坍缩（Neural collapse）是指在为分类而训练的深度神经网络中观察到的一种现象，其中学习到的特征和分类器权重在训练的末期表现出高度对称和结构化的行为 \cite{papyan2020prevalence,zhu2021geometric,yaras2022neural}。具体而言，在每个类别内部，特征坍缩到它们的类别均值；而在类别之间，这些均值变得最大程度地分离，形成一个单纯形等角配置（simplex equiangular configuration）。线性分类器在重新缩放的意义下与类别均值对齐。此外，最后一层分类器收敛于选择具有最近训练类别均值的类别。神经坍缩揭示了监督学习中出现的优化动力学、泛化和几何结构之间的深刻联系。
    \label{remark:neural-collapse}
\end{remark}

从上述分类示例可见，这样学习到的 $\vx$ 的表征 $\vz$ 实际上对应了一个非常简单的编码器：它几乎把每一类数据都映射到一个码字，即该类别的独热向量。从有损压缩角度看，这样的编码器损失过大，难以保留数据分布中的丰富信息。诸如图像分割、补全或生成等任务所需的其他信息，会在这种监督学习过程中严重丢失。

为弥补这一缺陷，我们希望学习一种不同的编码方案，使生成的表征 $\vz$ 能尽可能捕捉数据分布信息，而不仅限于对分类有用的信息。更准确地说，我们希望表征 $\vz = f(\vx, \theta)$ 满足：存在解码器 $g(\vz,\eta)$ 使
\begin{equation}
    \|\vx - \hat \vx\|_2 \le \epsilon \quad \mbox{s.t.} \quad \hat{\vx}
    = g(f(\vx)).
\end{equation}
我们将在第 \ref{ch:consistent-representations} 章中正式把这种编码-解码方案称为“自编码器”（autoencoder），并把这种表征 $\vz$ 称为“一致”（consistent）表征。此处可先将其视为上一节在原始数据 $\vx$ 域中引入的量化码 $\hat \vx$ 的推广。

信息瓶颈公式 \eqref{chap4-eqn:information-bottleneck} 并不要求学习到的表征 $\vz$ 是一致的。因此，至少应使用一个（一致的）表征学习（RL）目标对其进行修正：
\begin{eqnarray}
    \max_{\theta,\eta}\; \mbox{RL}(\vx, \y, \vz) \doteq I(\vz; \y) - \beta
    I(\vx; \vz) \nonumber \\ \mathrm{s.t.}\quad  \vz = f(\vx,\theta),\quad
    \|\vx - g(\vz, \eta) \|_2 \le \epsilon.
    \label{chap4-eqn:representation-learning}
\end{eqnarray}
在最一般的设置下求解上述问题非常艰巨。通常，我们甚至不知道哪类编码器或解码器能够保证一致性\footnote{在实践中，人们倾向于通过试错法、独立于上述目标来学习它们，正如我们将在第 \ref{ch:consistent-representations} 章中进一步讨论的那样。}。此外，由于表征 $\vz$ 的分布未知，计算互信息 $I(\vz;\y)$ 即使并非不可能，也极具挑战。

在后续章节讨论一般情形之前，本章余下部分将先展示：在数据分布由清晰（低维）类别组成的特殊设定下，如何在一定程度上解决该问题。

\subsection{线性判别表征}\label{subsec:LDR}

给定来自混合分布 $\mathcal{D}$ 的数据 $\vX$ 能否被有效分类或聚类，取决于分量分布 $\mathcal{D}_k$ 的可分离性（判别性）有多强（或可被提升到多强）。一个常见工作假设是：每个类别分布都具有相对 {\em 低维} 的内在结构。因此，可假设每个类别分布 $\mathcal{D}_k$ 的支撑位于低维子流形上，例如维度为 $d_k \ll D$ 的 $\mathcal{M}_k$；而 $\vx$ 的整体分布 $\mathcal D$ 的支撑位于高维环境空间 $\Re^D$ 中这些子流形的混合 $\mathcal M = \cup_{k=1}^K \mathcal{M}_k$ 上，如图 \ref{fig:representation-learning} 所示。

我们不仅需要识别低维分布，还希望以最有利于后续任务（如分类、聚类和条件生成，正如我们将来会看到的）的形式来表征该分布。为此，我们要求学习到的表征具有以下属性：
\begin{enumerate}
    \item {\em 类内可压缩（Within-Class Compressible）：} 来自同一类别的样本特征应该具有强 {\em 相关性}，即它们属于一个低维线性子空间。
    \item {\em 类间具判别性（Between-Class Discriminative）：} 来自不同类别的样本特征应该高度 {\em 不相关}，并且属于不同的不相干低维线性子空间。
    \item {\em 最大化多样性表征（Maximally Diverse Representation）：} 只要每个类别的特征与其他类别不相干，其维度（或方差）应该 {\em 尽可能大}。
\end{enumerate}
我们将这种表征称为 {\em 线性判别表征}（linear discriminative representation, LDR）。

请注意，第一个属性与我们在 \Cref{sub:pca} 中讨论的经典 {\em 主成分分析}（PCA）目标高度一致。其本质是：同一类别样本\footnote{包括同一样本的增强实例}的特征应当（或被变换为）高度相关。换言之，表征对相似样本应是 {\em 收缩的}（contractive）。

\begin{figure}
    \centering
    \includegraphics[width=0.8\linewidth]{\toplevelprefix/chapters/chapter4/figs/mcr_diagram.png}
    \caption{高维数据 $\vx\in \Re^D$ 的分布 $\mathcal D$ 支撑在流形 $\mathcal{M}$ 上，其类别支撑在低维子流形 $\mathcal{M}_k$ 上。我们的目标是学习一个由 $\theta$ 参数化的映射 $f(\vx, \theta)$，使得 $\vz_i = f(\vx_i, \theta)$ 位于最大程度不相关的子空间 $\{\mathcal{S}_k\}$ 的并集上。}
    \label{chap4-fig:mcr-diagram}
\end{figure}

第二个属性要求：对属于不同类别或被视为不相似的样本，其特征应被变换为（几何上）不相干或（统计上）独立，如图 \ref{fig:representation-learning} 所示。\Cref{chap4-fig:mcr-diagram} 则说明了数据分布是低维子流形混合时的情形。通过压缩（去噪或聚类），我们先识别出真实数据分布（左）；随后希望进一步变换该分布，使这些子流形最终成为彼此不相干/独立的线性子空间（右）。

\begin{remark}[线性判别分析]
    人们可能会注意到，第二个标准类似于经典的 {\em 线性判别分析}（LDA）~\cite{HastieTiFr09}。线性判别分析（LDA）~\cite{HastieTiFr09} 是一种监督降维技术，旨在找到最大化类别可分性的数据线性投影。具体而言，给定标记数据，LDA 寻求一种线性变换，将高维输入投影到较低维空间，在该空间中类别被最大程度地分离。请注意，PCA 是一种无监督方法，它将数据投影到最大方差的方向上，而不考虑类别标签。虽然 PCA 纯粹专注于保留全局方差结构，但 LDA 明确利用标签信息来增强判别能力；参见 \Cref{fig:LDA} 中的比较。
    \begin{figure}
        \centering
        \includegraphics[width=0.75\linewidth]{\toplevelprefix/chapters/chapter4/figs/LDA.png}\vspace{-0.1in}
        \caption{PCA 与 LDA 的比较。图片引自 \url{https://sebastianraschka.com/Articles/2014_python_lda.html}。}
        \label{fig:LDA}
    \end{figure}
    但 LDA 仅在原始数据上执行线性投影。正如我们将看到的，LDR 更一般：它寻求一种（可能非线性的）变换，使生成表征在线性结构上可分且彼此不相干。
\end{remark}

\begin{remark}[对比学习]
    人们可能会注意到，传统的 {\em 对比学习}（contrastive learning）\citep{hadsell2006dimensionality,Van_den_Oord2018-qf,he2019momentum} 旨在实现与前两个属性所规定的表征相似的目标。对于属于 $k$ 个不同类别的数据，如果 $k$ 很大，随机选择的一对 $(\vx_i, \vx_j)$ 属于不同类别的概率很高。\footnote{例如，当 $k \ge 100$ 时，随机一对属于不同类别的概率为 99\%。} 因此，理想情况下，样本 $\vx_i$ 的表征 $\vz_i = f(\vx_i, \bm \theta)$ 应该与其他样本 $\vx_j$ 的表征 $\vz_j$ 高度不相干，而与其变换版本 $\tau(\vx_i)$ 的特征（记为 $\vz(\tau(\vx_i))$，其中 $\tau$ 属于所考虑的某个增强集 $\mathcal T$）相干。因此，启发式地提出，为了提升学习到的表征的判别性，人们可以寻求最小化所谓的 {\em 对比损失}（contrastive loss）：
    \begin{equation}
        \min_{\bm \theta} - \log \frac{\exp(\langle \vz_i,
        \vz(\tau(\vx_i))\rangle)}{\sum_{j\not= i} \exp(\langle \vz_i,
        \vz_j\rangle)},
        \label{eqn:contrastive-loss}
    \end{equation}
    只要内积 $\langle \vz_i, \vz(\tau(\vx_i))\rangle$ 很大且对于 $i \not= j$ 的 $\langle \vz_i, \vz_j\rangle$ 很小，该损失就很小。这种对比损失及其许多变体已广泛用于流行的（自监督）特征学习方法中，例如用于视觉数据的 DINO 类型特征 \cite{caron2021unsupervisedlearningvisualfeatures,DINOv2}。正如我们将看到的，基于上一节研究的有损编码框架，我们将能够得出一个原则性的目标函数，该函数不仅促进了这两个属性，同时还显著简化了诸如学习 DINO 类型特征的实践（正如我们将在第 \ref{ch:applications} 章中详细介绍的那样）。
\end{remark}

第三个属性也很重要，因为我们希望学习到的特征能够揭示一个类别不同于所有其他类别的所有可能原因。例如，为了区分“苹果”和“橘子”，我们不仅关心颜色，还关心形状和叶子。理想情况下，每个子空间 $\{\mathcal{S}_k\}$ 的维度应该等于相应子流形 $\mathcal{M}_k$ 的维度。如果我们希望映射 $f(\vx,\theta)$ 对于图像生成等任务是 {\em 可逆的}，那么这个属性将非常重要。例如，如果我们从“苹果”的特征子空间中抽取不同的样本点，我们应该能够通过逆映射 $g(\vz, \eta)$ 对它们进行解码，以生成各种各样的苹果图像。这样一对互逆映射 $f$ 和 $g$ 构成了一个自编码（autoencoding），我们将在第 \ref{ch:consistent-representations} 章中详细研究。请注意，通过最小化交叉熵 \eqref{chap4-eqn:cross-entropy} 学习到的表征通常不具备此属性，因为它们会遭受前面讨论的所谓 {\em 神经坍缩} 现象（见注 \ref{remark:neural-collapse}）。

一般而言，尽管每个类别/聚类的内在结构可能是低维的，但它们在原始表征 $\vx$ 中通常并非简单线性（或高斯）结构，往往需要先经某种非线性变换后才“线性化”。\footnote{我们将在 \Cref{ch:consistent-representations} 中明确讨论如何实现这一点。} 因此，我们希望通过非线性变换 $f(\vx,\theta)$ 学到一种表征，使各类别对应的子空间尽可能不相干。更准确地说，我们希望学习映射 $\vz = f(\vx,\theta)$，将每个子流形 $\mathcal{M}_k \subset \Re^D$（\Cref{chap4-fig:mcr-diagram} 左）映射到一个 {\em 线性} 子空间 $\mathcal{S}_k \subset \Re^d$（\Cref{chap4-fig:mcr-diagram} 右）。在一定意义下，得到的多个子空间 $\{\mathcal{S}_k\}$ 可视为原始数据 $\bm x$ 所生成特征 $\vz$ 的具判别性的 {\em 广义主成分}（generalized principal components）\cite{GPCA}；若它们正交，则还可视为 {\em 独立成分}（independent components）\cite{hyvarinen2000independent}。正如下一章 \Cref{ch:deep-networks} 将展示的，深度网络恰好可用于建模并实现这种从数据分布到线性判别表征的非线性变换。

\subsection{最大编码率降低原理}\label{subsec:MCR2}

尽管线性判别表征（LDR）的三个属性——{\em 类间具判别性}、{\em 类内可压缩} 和 {\em 最大化多样性表征}——都是我们希望学习到的表征 $\vz$ 具备的性质，但要同时实现它们并不容易：这些属性是否彼此兼容，以至于可以一次性达成？如果可以，是否存在一个 {\em 简单且有原则} 的目标，用于综合衡量表征质量？这些问题的关键，在于为随机变量 $\vz$ 的分布（或其有限样本 $\{\bm z_i\}_{i=1}^N$）找到一个原则性的“紧凑性度量”或“信息增益”度量。该度量应能依据内在维度或体积，直接且准确地刻画分布的内在几何与统计属性。与交叉熵 \eqref{chap4-eqn:cross-entropy} 或信息瓶颈 \eqref{chap4-eqn:information-bottleneck} 不同，这种度量不应完全依赖类别标签，以便适用于更一般的场景，如监督、自监督、半监督与无监督学习（后续小节将进一步说明）。

不失一般性，设随机向量 $\vx$ 的分布 $\mathcal D$ 是若干分量分布的混合，即 $\mathcal D = \cup_{k=1}^K \mathcal{D}_k$，其中每个 $\mathcal{D}_k \subset \Re^D$ 在高维环境空间 $\Re^D$ 中具有较低内在维度。令 $\bm X_k \in \Re^{D\times N_k}$ 表示数据矩阵，其列向量是从分布 $\mathcal{D}_k$ 中抽取的样本，$N_k$ 为第 $k$ 类样本数。记 $\bm X=[\bm X_1,\dots,\bm X_K] \in \Re^{D\times N}$ 为全部样本，其中 $N=\sum_{k=1}^K N_k$。同时我们也记 $\vX=[\vx_1,\dots,\vx_N]$，其中 $\vx_i$ 是第 $i$ 个样本。在编码映射
\begin{equation}
    \vx  \xrightarrow{\hspace{2mm} f(\vx)\hspace{2mm}} \vz,
\end{equation}
下，对每个 $i=1,\dots,N$，输入样本映射为 $\vz_i = f(\vx_i)$。稍作符号滥用，我们也写作 $\vZ_k = f(\vX_k)$ 和 $\vZ = f(\vX)$。因此，既有 $\vZ = [\bm Z_1,\dots,\bm Z_K]$，也有 $\vZ = [\vz_1,\dots,\vz_N]$。

一方面，为使学习到的特征具有判别性，不同类别/聚类的特征应尽可能彼此 {\em 不相干}。因此，它们整体应张成尽可能大体积（或高维度）的空间，整个集合 $\vZ$ 的编码率也应尽可能大。另一方面，同一类别/聚类内的特征应高度相关且相干，因此每一类应仅张成体积较小的空间（或子空间），其编码率应尽可能小。下面我们介绍如何度量学习特征的编码率。

\paragraph{特征的编码率。} 需要注意的是，评估编码率的一个实际困难在于：表征 $\vZ$ 的潜在分布通常未知。为此，我们可将特征 $\vZ = [\vz_1, \ldots, \vz_N]$ 近似视为来自多元高斯分布的样本。在该假设下，正如 \Cref{subsec:lossy DR} 所述，特征 $\vZ$ {\em 作为整体} 的紧凑性可由每个样本的平均编码长度来衡量，即 {\em 编码率}（coding rate）。给定精度 $\epsilon > 0$（见 \eqref{eqn:rate-Gaussian}），其定义为
\begin{equation}
    R_{\epsilon}(\vZ) = \frac{1}{2}\log\det\left(\I +
    \frac{d}{N\epsilon^{2}}\vZ\vZ^{\top}\right).
    \label{chap4-eqn:coding-length-eval}
\end{equation}

另一方面，我们希望非线性变换 $f(\vx)$ 将各类别子流形 $\mathcal{M}_k \subset \R^D$ 映射到尽可能不相干的线性子空间 $\mathcal{S}_k \subset \R^d$，使学习到的特征 $\vZ$ 落在若干低维子空间的并集中。该结构使我们能够分别分析各子空间，从而更准确地评估编码率。%
回想一下，$\vZ_k$ 的列向量表示 $\vX_k$（$k=1,\dots,K$）中样本的特征。于是，$\bm Z_k$ 中特征的编码率可写为
\begin{align}
    R_{\epsilon}(\vZ_k) = \frac{N_k}{2N}\log\det\left(\I +
    \frac{d}{N_k\epsilon^{2}}\vZ_k\vZ_k^{\top}\right)
\end{align}
那么，每个类别中特征的平均编码率之和为
\begin{equation}
    R_{\epsilon}^c(\vZ) \doteq \sum_{k=1}^K R_{\epsilon}(\vZ_k),
    \label{chap4-eqn:compress-loss-eval}
\end{equation}

因此，$\vX$ 的一个好的表征 $\vZ$ 是能够在整体编码率和所有类别的编码率之间实现巨大差异的表征：
\begin{equation}
    \Delta R_{\epsilon}(\vZ) \doteq R_{\epsilon}(\vZ) - R_{\epsilon}^c(\vZ).
    \label{chap4-eqn:coding-length-reduction}
\end{equation}
根据本章前文讨论，这一差异可解释为：在整个集合 $\vZ$ 中识别出正确低维聚类 $\vZ_k$ 后所获得的“信息增益”。

若将特征映射 $f(\cdot)$ 选为带参数 $\theta$ 的深度神经网络 $f(\cdot,\theta)$，则表征过程及由此产生的率降低可写为
\begin{equation}
    \vX
    \xrightarrow{\hspace{2mm} f(\vx, \theta)\hspace{2mm}} \vZ
    \xrightarrow{\hspace{2mm} \epsilon \hspace{2mm}} \Delta
    R_{\epsilon}(\vZ).
    \label{chap4-eqn:flow}
\end{equation}
请注意，$\Delta R_{\epsilon}$ 对特征 $\vZ$ 的尺度是 {\em 单调} 的。为保证不同表征之间可公平比较，必须对学习特征 {\em 做尺度归一化}。一种做法是约束每个类别 $\vZ_k$ 的 Frobenius 范数与其特征数成比例，即对 $\vZ_k \in \mathbb R^{d \times N_k}$ 令 $\|\vZ_k\|_F^2 = N_k$；另一种做法是将每个特征归一化到单位球面，即 $\vz_i \in \mathbb{S}^{d-1}$，其中 $N_k$ 为第 $k$ 类样本数。这一表述也为深度网络训练中“批量归一化”（batch normalization）的必要性提供了自然解释 \cite{ioffe2015batch}。%

当表征具有可比性后，目标就变为学习特征 $\vZ  = f(\vX, \theta)$，使“整体编码率”与“分类别编码率之和”之间的降低量最大：
\begin{equation}
    \begin{aligned}
        \max_{\theta } & \;  \Delta R_{\epsilon}\big(\vZ \big) \doteq
        R_{\epsilon}(\vZ) - R_{\epsilon}^c(\vZ ), \\
        \mbox{s.t.} & \ \ \, \vZ = f(\bm X, \theta),\  \|\vZ_k\|_F^2
        = N_k,\ k=1,\dots,K.
    \end{aligned}
    \label{eqn:maximal-rate-reduction}
\end{equation}
我们将其称为 {\em 最大编码率降低}（maximal coding rate reduction, MCR$^2$）原理，它体现了亚里士多德名言：
\begin{quote}
    \centering
    “{\em 整体大于部分之和。}”
\end{quote}
为了学习最佳表征，我们要求 {\em 整体尽可能大于部分之和}。在很大程度上，量 $\Delta R$ 可视为目标 \eqref{chap4-eqn:representation-learning} 中互信息项 $I(\vz; \y)$ 的一种可计算实现。

\begin{example}
    为了进一步阐明这些想法，让我们检查 \Cref{fig:informative-representation} 中所示的例子。
    \begin{figure}[t]
        \centering
        \includegraphics[width=0.99\linewidth]{\toplevelprefix/chapters/chapter4/figs/Expansion2.png}
        \caption{在识别出低维数据分布之后，我们希望进一步将数据分布转换为更具信息量的结构化表征：$R$ 是覆盖整个空间的 $\epsilon$-球的数量，$R^c$ 是所有子空间（绿球）的数量之和。$\Delta R$ 是它们的差值（蓝球的数量）。
        }\label{fig:sphere-packing}
        \label{fig:informative-representation}
    \end{figure}
    从压缩的角度来看，右侧的表征是 {\em 最紧凑的}，其意义在于，当所有特征被编码为单个高斯分布（蓝色）时的编码率，与特征被正确聚类并编码为两个独立子空间（绿色）时的编码率之间的差异是最大的。\footnote{直观地说，所有特征张成的整个空间的“体积”与特征实际占据的体积之间的比率是最大的。}
\end{example}

请注意，上述 MCR$^2$ 原理最初是为监督学习设计的，即组隶属度（或类别标签）已知。通过引入隶属度矩阵，该原理可自然扩展到无监督学习：该矩阵编码每个数据点到潜在组/聚类的分配（可为软分配）。具体地，令 $\bm \Pi = \{\bm \Pi_k\}_{k=1}^K \subset \R^{N\times N}$ 为一组对角矩阵，其对角元素表示 $N$ 个样本到 $K$ 个类别的隶属度，即 $\bm \Pi$ 属于单纯形 $\Omega \doteq \{\bm \Pi: \bm \Pi_k \ge \bm 0: \sum_{k=1}^K \bm \Pi_k = \bm I_N\}$。据此可定义关于划分 $\bm \Pi$ 的平均编码率为
\begin{align}\label{eq:MCRc}
    R_{\epsilon}^c(\vZ \mid \bm \Pi) \doteq \sum_{k=1}^K
    \frac{\mathrm{tr}(\bm \Pi_k)}{2N}\log\det\left(\bm I +
        \frac{d}{\mathrm{tr}(\bm \Pi_k)\epsilon^2}\bm Z\bm \Pi_k \bm
    Z^\top \right).
\end{align}
当 $\bm Z$ 给定时，$R_{\epsilon}^c(\vZ \vert \bm \Pi)$ 关于 $\bm \Pi$ 是凹函数。于是，无监督学习下的 MCR$^2$ 原理可写为
\begin{align}\label{eq:MCR pi}
    \max_{\bm \Pi, \theta} & \  \Delta R_{\epsilon}\big(\vZ  \mid \bm \Pi)
    \doteq R_{\epsilon}(\bm Z) - R_{\epsilon}^c(\vZ \mid \bm \Pi) \notag \\
    \mathrm{s.t.}  & \ \ \ \bm Z = f(\bm X, \theta),\ \|\bm Z \bm
    \Pi_k\|_F^2 = N_k,\ k = 1,\dots,K,\ \bm\Pi \in \Omega.
\end{align}
与 \eqref{eqn:maximal-rate-reduction} 相比，上式允许联合优化组隶属度与网络参数。特别地，当 $\bm \Pi$ 固定为将 $N$ 个数据点分配到 $K$ 个组的隶属度矩阵时，问题 \eqref{eq:MCR pi} 就退化为问题 \eqref{eqn:maximal-rate-reduction}。

\begin{figure}[t]
    \centering
    \includegraphics[width=0.9\linewidth]{\toplevelprefix/chapters/chapter4/figs/mcr2-global.png}
    \caption{{\bf 局部优化景观：} 根据 \Cref{thm:MCR2-properties}，率降低目标的全局最大值对应于具有相互不相干子空间的解。}
    \label{fig:mcr-global}
\end{figure}

\begin{remark}[计算编码率]
    一个经典事实是，计算一个 \(n \times n\) 矩阵的对数行列式需要 \(\cO(n^{3})\) 的时间。因此，在现代数据分析的规模下，计算编码率 \(\Delta R_{\epsilon}\) 或 \(R_{\epsilon}\) 等量似乎并不可行。对此有几种处理方法；在本书后续实验中，我们会直接使用这些方法而不再逐一说明。
    \begin{itemize}
        \item \textit{利用编码率的结构。} 假设 \(\vZ \in \R^{d \times n}\)。则有
            \begin{equation}
                \logdet(\vI + \alpha \vZ\vZ^{\top}) = \logdet(\vI +
                \alpha \vZ^{\top}\vZ).
            \end{equation}
            因此，计算矩阵乘积并进一步计算编码率所需的时间量级为 \(\min\{d^{2}n + d^{3}, dn^{2} + n^{3}\}\)。
        \item \textit{利用 Gram 矩阵的对称性。}
            注意，\(\vZ\vZ^{\top}\)（以及类似的 \(\vZ^{\top}\vZ\)）是对称矩阵。因此，计算其特征值 \(\lambda_{i}(\vZ\vZ^{\top})\) 相对便宜，尽管仍然是三次时间复杂度。此时，对数行列式可以写为
            \begin{equation}
                \logdet(\vI + \alpha \vZ\vZ^{\top}) = \sum_{i =
                1}^{d}\log(1 + \alpha \lambda_{i}(\vZ\vZ^{\top})),
            \end{equation}
            对 \(\vZ^{\top}\vZ\) 也同理。不过，下一种方法利用了输入矩阵的更多结构，通常效果更好，因此这种方法在实践中较少使用。
        \item \textit{利用对数行列式输入矩阵的正定性。}
            由于 \(\vZ\vZ^{\top}\)（以及类似的 \(\vZ^{\top}\vZ\)）是对称半正定矩阵，\(\vI + \alpha \vZ\vZ^{\top}\) 是对称正定矩阵。这意味着我们可以计算 Cholesky 分解 \(\vL\vL^{\top} = \vI + \alpha \vZ\vZ^{\top}\)，其中 \(\vL \in \R^{d \times d}\) 是对角元素为正的下三角矩阵。于是
            \begin{equation}
                \logdet(\vI + \alpha \vZ\vZ^{\top}) = 2\sum_{i =
                1}^{d}\log(L_{ii}).
            \end{equation}
            这种方法在实践中通常最便宜，尤其是在同时根据规模选择计算 \(\vZ\vZ^{\top}\) 还是 \(\vZ^{\top}\vZ\) 时。不过，它仍然保持三次时间复杂度。
    \end{itemize}
\end{remark}

\vspace{-0.1in}

\subsection{编码率降低的优化属性}
\label{sec:MCR-landscape}

在本小节中，我们通过分析 MCR$^2$ 函数的最优解及其优化景观的结构来研究其优化属性。为了避开神经网络引入的技术困难，我们考虑问题 \eqref{eqn:maximal-rate-reduction} 的简化版本如下：
\begin{align}\label{eqn1:maximal-rate-reduction}
    \max_{\bm Z}\ R_{\epsilon}(\vZ) - R_{\epsilon}^c(\vZ)\qquad
    \mathrm{s.t.}\quad \|\bm Z_k\|_F^2 = N_k,\ k =1,\dots,K.
\end{align}
理论上，MCR$^2$ 原理 \eqref{eqn1:maximal-rate-reduction} 具有很强的一般性，并且可以应用于 {\em 任何} 分布的表征 $\vZ$，只要这些分布的率 $R_\epsilon$ 和 $R^c_\epsilon$ 能够被准确且高效地评估。最优表征 $\vZ^{\ast}$ 应该具有一些有趣的几何和统计属性。我们在这里通过子空间的特例揭示了最优表征的良好属性，子空间在机器学习中有许多重要的用例。当 $\vZ$ 的期望表征是多个子空间时，\eqref{eqn1:maximal-rate-reduction} 中的率 $R_\epsilon$ 和 $R^c_\epsilon$ 分别由 \eqref{chap4-eqn:coding-length-eval} 和 \eqref{chap4-eqn:compress-loss-eval} 给出。在最大率降低时，MCR$^2$ 实现了其最优表征，记为 $\vZ^{\ast} = [\vZ_1^*,\dots,\vZ_K^*]$，其中 $\rank{(\vZ_{k}^*)}\le d_k$。可以证明 $\vZ^{\ast}$ 具有以下期望的属性（正式陈述和详细证明见 \cite{yu2020learning}）。

\begin{theorem}[\bf 全局最优解的刻画]
    假设 $\vZ^{\ast} = [\vZ_1^*,\dots,\vZ_K^*]$ 是问题~\eqref{eqn1:maximal-rate-reduction} 的全局最优解。以下陈述成立：
    \begin{itemize}
        \item {\em 类间具判别性}：只要环境空间足够大（$d \ge \sum_{k=1}^{K} d_k$），子空间彼此之间都是正交的，{\em 即}，对于 $k \not= l$，有 $(\vZ_{k}^*)^{\top} \vZ_{l}^* = \bm{0}$。
        \item {\em 最大化多样性表征}：
            只要编码精度足够高，即 $\epsilon ^4 < c\cdot \min_{k}\left\{ \frac{N_k}{N}\frac{d^2}{d_k^2}\right\}$，其中 $c>0$ 是一个常数。每个子空间都达到其最大维度，即 $\mathrm{rank}{(\vZ_{k}^*)}= d_k$。此外，$\vZ_{k}^*$ 的最大的 $d_k-1$ 个奇异值相等。
            \label{thm:MCR2-properties}
    \end{itemize}
\end{theorem}

该定理表明，MCR$^2$ 原理促进了将数据嵌入到多个独立的子空间中（如 \Cref{fig:mcr-global} 所示），并且特征在每个子空间中呈 {\em 各向同性} 分布（可能除了一个维度之外）。值得注意的是，该定理还证实了通过 MCR$^2$ 原理学习到的特征表现出 \Cref{subsec:LDR} 中讨论的期望的低维判别属性。此外，在所有此类判别表征中，它偏好在环境空间中具有最高维度的表征。这与信息瓶颈~\eqref{chap4-eqn:information-bottleneck} 的目标有本质的不同。

\paragraph{正则化 MCR$^2$。}
上述定理刻画了率降低目标的全局最优解的属性。那么其他最优解（例如局部最优解）呢？由于 Frobenius 范数的约束，从优化理论的角度分析问题 \eqref{eqn1:maximal-rate-reduction} 是一项困难的任务。因此，我们考虑 \eqref{eqn1:maximal-rate-reduction} 的拉格朗日公式。这可以被视为 \eqref{eqn1:maximal-rate-reduction} 的紧松弛甚至等价问题，在正则化参数的特定设置下，它们的最优解是一致的；参见 \cite[Proposition 3.2]{wang2024global}。具体而言，我们研究的公式（此后称为 \textit{正则化 MCR$^2$ 问题}）如下：
\begin{align}\label{eq:MCR-reg}
    \max_{\bm Z}\ R_{\epsilon}(\vZ) - R_{\epsilon}^c(\vZ) -
    \frac{\lambda}{2}\|\bm Z\|_F^2,
\end{align}
其中 $\lambda > 0$ 是正则化参数。尽管规划 \eqref{eq:MCR-reg} 是高度非凹的，并且在其梯度计算中涉及矩阵求逆，我们仍然可以明确刻画其局部和全局最优解如下（详细证明见 \cite[Theorem 3.1]{wang2024global}）。

\begin{theorem}[\bf 局部和全局最优解]\label{thm:mcr-global-opt}
    令 $N_k$ 表示每个 $k \in \{1,\dots,K\}$ 的第 $k$ 个类别中的训练样本数，$N_{\max} \doteq \max\{N_1,\dots,N_K\}$，$\alpha=d/(N\epsilon^2)$，以及对于每个 $k \in \{1,\dots,K\}$ 的 $\alpha_{k} = d/(N_k\epsilon^2)$。给定编码精度 $\epsilon > 0$，如果正则化参数满足
    \begin{align}\label{eq:lambda}
        \lambda \in \left(0,
            \frac{d(\sqrt{N/N_{\max}}-1)}{N(\sqrt{N/N_{\max}}+1)\epsilon^2}
        \right],
    \end{align}
    则以下陈述成立：\\
    (i) ({\bf 局部最大值点}) $\bm Z^* = \left[\bm Z_1^*,\dots,\bm Z_K^* \right]$ 是问题 \eqref{eq:MCR-reg} 的局部最大值点，当且仅当第 $k$ 个块允许以下分解
    \begin{align}\label{eq:Zk opti}
        \bm Z_k^* = \left(\frac{ \eta_k + \sqrt{\eta_k^2 -
        4\lambda^2N/N_k}}{2\lambda \alpha_{k}}\right)^{1/2} \bm
        U_k \bm V_k^\top,
    \end{align}
    其中 (a) $r_k = \mathrm{rank}(\bm Z_k^*)$ 满足 $r_k \in [0,\min\{N_k,d\})$ 且 $\sum_{k=1}^K r_k \le \min\{N,d\}$，(b) $\bm U_k \in \mathcal{O}^{d \times r_k}$ 满足对于所有 $k \neq l$ 有 $\bm U_k^{\top}\bm U_l = \bm 0$，$\bm V_k \in \mathcal{O}^{N_k \times r_k}$，以及 (c) 对于每个 $k\in \{1,\dots,K\}$，$\eta_k := (\alpha_k-\alpha) - \lambda(N/N_k+1)$。
    \\
    (ii) ({\bf 全局最大值点}) $\bm Z^* = \left[\bm Z_1^*,\dots,\bm Z_K^* \right]$ 是问题 \eqref{eq:MCR-reg} 的全局最大值点，当且仅当 (a) 它满足上述所有条件且 $\sum_{k=1}^K r_k = \min\{N,d\}$，以及 (b) 对于所有满足 $N_k < N_l$ 且 $r_l > 0$ 的 $k \neq l \in [K]$，我们有 $r_k = \min\{N_k,d\}$。
\end{theorem}

\begin{figure}[t]
    \centering
    \includegraphics[width=0.9\linewidth]{\toplevelprefix/chapters/chapter4/figs/mcr2-global-local.png}
    \caption{{\bf 全局优化景观：} 根据 \cite{sun2015nonconvex,lee2016gradient}、\Cref{thm:mcr-global-opt,thm:mcr-benign-opt-landscape}，（正则化）率降低目标的全局和局部最大值都对应于具有相互不相干子空间的解。所有其他临界点都是严格鞍点。}
    \label{fig:mcr-global-local}
\end{figure}

该定理明确刻画了问题 \eqref{eq:MCR-reg} 的局部和全局最优解。直观地说，这表明问题 \eqref{eq:MCR-reg} 的每个局部最大值点所对应的特征都是低维且具判别性的。尽管我们在定理~\ref{thm:mcr-global-opt} 中刻画了局部和全局最优解，但仍然未知是否可以使用梯度下降（GD）高效地计算这些解来解决问题 \eqref{eq:MCR-reg}，因为 GD 可能会卡在其他临界点（如鞍点）上。幸运的是，\cite{sun2015nonconvex,lee2016gradient} 表明，如果一个函数是二次连续可微的并且满足 {\em 严格鞍点属性}（strict saddle property），即每个临界点要么是局部最大值点，要么是严格鞍点\footnote{如果问题 \eqref{eq:MCR-reg} 的临界点具有严格正曲率的方向~\cite{sun2015nonconvex}，我们称其为严格鞍点。这包括具有严格正曲率的经典鞍点以及局部最小值点。}，那么在随机初始化的情况下，GD 几乎必然收敛到其局部最大值点。我们通过刻画问题 \eqref{eq:MCR-reg} 的所有临界点来研究其全局优化景观如下（详细证明见 \cite[Theorem 3.3]{wang2024global}）。

\begin{theorem}[\bf 良性全局优化景观]\label{thm:mcr-benign-opt-landscape}
    给定编码精度 $\epsilon > 0$，如果正则化参数满足 \eqref{eq:lambda}，则问题 \eqref{eq:MCR-reg} 的任何临界点 $\bm Z$ 要么是局部最大值点，要么是严格鞍点。
\end{theorem}
综合起来，上述两个定理表明，与率降低目标的每个局部最大值点（不仅是全局最大值点）相关联的学习特征都被结构化为不相干的低维子空间。此外，（正则化）率降低目标 \eqref{eqn:maximal-rate-reduction} 具有非常良性的景观，仅有局部最大值点和严格鞍点作为临界点，如 \Cref{fig:mcr-global-local} 所示。根据 \cite{sun2015nonconvex,lee2016gradient}，\Cref{thm:mcr-global-opt,thm:mcr-benign-opt-landscape} 意味着通过从随机初始化开始对率降低目标 \eqref{eqn:maximal-rate-reduction} 应用（随机）梯度下降，可以高效地找到低维且具判别性的表征（LDR）。这些结果也间接解释了为什么在 \Cref{eg:Rate-Reduction-CIFAR10} 中，如果选择的网络具有足够的表达能力并且训练良好，生成的表征通常会给出一个不相干的线性表征，这很可能对应于全局最优解。感兴趣的读者可以参考 \cite{wang2024global} 获取证明。

\section{学习图像数据的表征}
\label{sec:representations-image}
\subsection{监督表征学习}
\label{sec:MCR2-classification}
\label{eg:Rate-Reduction-CIFAR10}

我们在这里展示 MCR$^2$ 目标 \eqref{eqn:maximal-rate-reduction} 如何帮助学习比交叉熵（CE）\eqref{chap4-eqn:cross-entropy} 更好的图像分类表征。在这里，为了简单起见，\footnote{我们将在下一章 \Cref{ch:deep-networks} 中严格研究并论证如何从 MCR$^2$ 目标出发并为其设计或推导合适的深度架构。} 我们采用流行的神经网络架构 ResNet-18~\cite{he2016deep} 来对特征映射 $\vz = f(\vx,\theta)$ 进行建模。我们通过反向传播算法（详见附录 \ref{app:BP-section}）优化神经网络参数 $\theta$，以最大化编码率降低。关于此实验的更多实现细节可以在 \cite{chan2021redunet} 中找到。

\Cref{fig:train-test-loss} 说明了两个编码率及其差值（在训练和测试数据上）随训练轮次的演变：在初始阶段之后，$R_\epsilon$ 逐渐增加而 $R^c_\epsilon$ 减少，这表明特征 $\bm Z$ 作为整体正在扩展，而每个类别 $\bm Z_k$ 正在被压缩。

我们首先使用 CIFAR10 图像分类数据集~\cite{krizhevsky2009learning} 评估性能，结果见 \Cref{fig:CIFAR10-performance} 中的表格。可以看到，MCR$^2$ 目标实现了与交叉熵（CE）损失相当的性能。此外，MCR$^2$ 目标对类别标签扰动更具鲁棒性，因为它为每个类别学习的是联合（子空间）表征，而不是拟合单个样本的标签。

\begin{figure}
    \centering
    \includegraphics[width=0.6\textwidth]{\toplevelprefix/chapters/chapter4/figs/loss_log.pdf}

    \caption{训练过程中 $R_\epsilon$、$R^c_\epsilon$ 和 $\Delta R_\epsilon$ 的演变。}
    \label{fig:train-test-loss}
\end{figure}

\begin{figure}
    \centering
    \includegraphics[width=\linewidth]{chapters/chapter4/figs/Table-label-corruption.png}
    \caption{在不同比例的类别标签被随机破坏的情况下，使用交叉熵（CE）损失或 MCR$^2$ 目标进行训练的分类准确率比较。}
    \label{fig:CIFAR10-performance}
\end{figure}

\begin{figure}[t]
    \begin{subfigure}[t]{0.45\textwidth}
        \centering
        \includegraphics[width=\textwidth]{\toplevelprefix/chapters/chapter4/figs/mcr_pca_individual.pdf}
        \caption{MCR$^2$ 学习到的各个类别特征的主成分分析（PCA）。}
        \label{fig:MCR-pca}
    \end{subfigure}
    \hfill
    \begin{subfigure}[t]{0.45\textwidth}
        \centering
        \includegraphics[width=\textwidth]{\toplevelprefix/chapters/chapter4/figs/ce_pca_individual.pdf}
        \caption{交叉熵学习到的各个类别特征的主成分分析（PCA）。}
        \label{fig:CE-pca}
    \end{subfigure}

    \caption{\small MCR$^2$ 与交叉熵学习到的特征的主成分比较。}
    \label{fig:feature-pca}
\end{figure}

\begin{figure*}[t]
    \begin{center}

        \includegraphics[width=0.45\textwidth]{\toplevelprefix/chapters/chapter4/figs/heatmap_mcr2.pdf}
        \hspace{0.25cm}
        \hfill
        \includegraphics[width=0.45\textwidth]{\toplevelprefix/chapters/chapter4/figs/heatmap_ce.png}
        \caption{\small 使用 MCR$^2$ 目标（\textbf{左}）和 CE 损失（\textbf{右}）学习到的特征之间的余弦相似度。}
        \label{fig:heatmap-plot}
    \end{center}
    \vskip -0.1in
\end{figure*}

\Cref{fig:feature-pca} 展示了每个类别学习到的特征 $\vZ_k$ 的奇异值分布。从 \Cref{fig:MCR-pca} 中显示的奇异值分布可以看出，MCR$^2$ 为每个类别学习到的特征 $\vz$ 分布在维度约为 10-12 的子空间上，而交叉熵 \Cref{fig:CE-pca} 学习到的特征维度要低得多。\footnote{如果继续增加迭代次数进行训练，CE 特征甚至可能坍缩到单一维度，这就是前面在 \Cref{remark:neural-collapse} 中讨论过的神经坍缩（neural collapse）现象。}

\Cref{fig:heatmap-plot} 展示了按类别排序的学习特征之间的余弦相似度。我们还比较了使用交叉熵（CE）\eqref{chap4-eqn:cross-entropy} 和 MCR$^2$ 目标 \eqref{eqn:maximal-rate-reduction} 学习到的特征的相似度。从图中可以清楚地看到，与使用交叉熵损失学习到的特征相比，MCR$^2$ 目标为不同类别学习到的特征更加不相干（或独立）。在每个类别内部，特征的相干性也较低，因此更加多样化，正如 \Cref{fig:MCR-pca} 中其奇异值分布所表明的那样。

我们将每个类别特征子空间中沿不同奇异向量映射的图像进行了可视化，我们称之为“主（principal）”图像。图 \ref{fig:pca-images} 展示了沿前 10 个奇异向量的这些主图像。正如我们所见，沿同一个奇异向量的图像在外观和结构上非常相似，而沿不同奇异向量的图像则彼此截然不同。
\begin{figure}
    \centering
    \includegraphics[width=\linewidth]{chapters/chapter4/figs/visual_2class.pdf}
    \caption{CIFAR10 图像数据集中“鸟（Bird）”和“船（Ship）”类别的前 10 个“主”图像。}
    \label{fig:pca-images}
\end{figure}

尽管 MCR$^2$ 目标似乎能学习到更好的表征，但上述实验所用网络架构仍缺乏明确的理论依据。目前尚不清楚为何这里采用的网络（ResNet-18）适合作为表征映射 $f(\vx, \theta)$，更不用说解释其内部学习到的层算子和参数 $\theta$。在下一章 \ref{ch:deep-networks} 中，{\em 我们将展示如何以“白盒”方式，从期望目标（如编码率降低）出发推导网络架构与组件}。

\subsection{无监督表征学习}
\label{sec:SimDINO}
% \yima{Ziyang，也许在这里介绍 SimDINO 的目标？并给出一些实验证明，将实现细节留到第 7 章（或新的第 8 章）。}

正如我们在上一节中所看到的，类别信息使得学习（混合）数据分布的结构化表征变得相对容易。然而，这种信息通常是不可用的，至少在海量规模上是如此。

现在我们介绍如何将 MCR$^2$ 原理扩展到无监督设置。与已知类别标签的监督学习相反，在无监督学习中，数据样本的组别归属是未知的。在这种情况下，我们可以对每个样本应用某类数据增强或变换，例如服从分布 $P_{\tau}$ 的 $\tau$。例如，对于图像，数据增强通常包括对每个图像样本进行平移、旋转或裁剪。我们可以将每个数据点的增强样本视为属于同一个组或簇。

更具体地说，给定一组 $N$ 个数据样本 $\vX =
[\vx_1,\dots,\vx_N]$，我们为每个数据点 $\vx_k$ 生成 $m$ 个增强视图
$\{\tau_i\}_{i=1}^m$，记为
$\tilde{\vX}_{k} = [\tau_1(\vx_k), \tau_2(\vx_k), \dots, \tau_m(\vx_k)]$。
我们将增强样本 $\tilde{\vX}_{k}$ 视为第 $k$ 个类别，总共有 $K=N$ 个组/簇。现在，我们可以像之前在监督学习情况中那样应用 MCR$^2$ 原理。具体而言，我们通过最大化编码率降低来学习特征映射 $f(\vx,\theta)$，如下所示：
\begin{equation}
    \begin{aligned}
        \max_{\theta } & \;  \Delta R_{\epsilon}\big(\vZ \big) \doteq
        R_{\epsilon}(\vZ) - R_{\epsilon}^c(\vZ ), \\
        \mbox{s.t.} & \ \ \, \vZ = f(\tilde{\vX}, \theta),\  \|\vZ_k\|_F^2
        = m,\ k=1,\dots,N.
    \end{aligned}
    \label{eqn:maximal-rate-reduction-unsupervised}
\end{equation}
这里，$\tilde{\vX} = [\tilde{\vX}_1,\dots,\tilde{\vX}_N]$，且 $\vZ_k =
f(\tilde{\vX}_k,\theta)$ 表示每个 $k=1,\dots,N$ 的 $\vx_k$ 的增强样本的特征。

让我们再次采用 ResNet-18 作为特征映射 $f(\vx,\theta)$，并优化网络参数 $\theta$ 以最大化无监督目标 \eqref{eqn:maximal-rate-reduction-unsupervised}。在 CIFAR10 数据集上训练后，我们在 \Cref{fig:train-test-loss-selfsup-mcr} 中可视化了两个编码率及其差值（针对训练和测试数据）随训练轮数（epochs）的演变。与监督学习的情况类似，在初始阶段之后，$R_\epsilon$ 逐渐增加而 $R^c_\epsilon$ 减小，这表明特征 $\bm Z$ 作为一个整体正在扩张，而每个簇 $\bm Z_k$ 正在被压缩。与 \Cref{fig:train-test-loss} 中显示的监督情况相比，我们可以观察到，在无监督情况下，总编码率 $R_\epsilon$ 迅速扩张，整体 MCR$^2$ 目标在局部最优处饱和。为了缓解这个问题，我们可以通过控制 $R_\epsilon$ 和 $R^c_\epsilon$ 的动态来改变优化策略。具体而言，我们可以通过将 \eqref{eqn:maximal-rate-reduction-unsupervised} 中的编码率 $R_\epsilon(\vZ)$ 替换为
$$\tilde{R}_\epsilon(\vZ) \doteq \frac{1}{2\gamma_1} \log \det
\left(\I + \frac{\gamma_2 d}{N m \epsilon^{2}}\vZ \vZ ^{\top}\right).$$
来重新缩放编码率。通过设置 $\gamma_1 = \gamma_2 = N$，我们得到了一个修改后的学习目标，称为 MCR$^2$-{\scriptsize CTRL}，其训练动态如 \Cref{fig:train-test-loss-selfsup-mcr-ctrl} 所示。正如我们所见，使用这个受控的 MCR$^2$ 目标，特征首先被压缩，然后逐渐扩张，从而导致整体目标 $\Delta R_{\epsilon}(\vZ)$ 达到更高的值。

为了评估学习到的表征的质量并验证它们是否具有良好的子空间结构，我们采用了子空间聚类算法 EnSC \cite{you2016oracle}。对于评估指标，我们考虑归一化互信息（NMI）、聚类准确率（ACC）和调整兰德指数（ARI）\footnote{有关它们的具体定义，请参阅 \cite{yu2020learning} 中的附录。}。为了进行比较，我们考虑了包括 JULE \cite{yang2016joint}、RTM \cite{nina2019decoder}、DEC \cite{xie2016unsupervised}、DAC \cite{chang2017deep} 和 DCCM \cite{wu2019deep} 在内的方法。聚类结果总结在表 \ref{fig:unsupervised-clustering-results} 中。正如我们所见，通过最大化 MCR$^2$-{\scriptsize CTRL} 目标学习到的特征比其他算法学习到的特征取得了更好的聚类准确率。
\begin{figure}[t]
    \begin{subfigure}[t]{0.45\textwidth}
        \centering
        \includegraphics[width=\textwidth]{\toplevelprefix/chapters/chapter4/figs/loss_selfsup_mcr.png}
        \caption{MCR$^2$}
        \label{fig:train-test-loss-selfsup-mcr}
    \end{subfigure}
    \hfill
    \begin{subfigure}[t]{0.45\textwidth}
        \centering
        \includegraphics[width=\textwidth]{\toplevelprefix/chapters/chapter4/figs/loss_selfsup_mcr_ctrl.png}
        \caption{MCR$^2$-ctrl}
        \label{fig:train-test-loss-selfsup-mcr-ctrl}
    \end{subfigure}
    \caption{\small 在 CIFAR-10 上进行无监督学习的训练过程中 MCR$^2$ 编码率的演变。}
    \label{fig:train-test-loss-selfsup}
\end{figure}

\begin{figure}
    \centering
    \includegraphics[width=\linewidth]{chapters/chapter4/figs/clustering_results.png}
    \caption{在 CIFAR10、CIFAR100 和 STL10 数据集上，使用 MCR$^2$-{\scriptsize CTRL} 目标与其他无监督学习算法进行训练的聚类准确率比较。}
    \label{fig:unsupervised-clustering-results}
\end{figure}

然而，这里有一个主要的缺点：基于 MCR$^2$ 的无监督学习方法在计算 $R^c_\epsilon$ 时，\textit{每个样本}都需要进行复杂度为 $\O(d^3)$ 的昂贵 $\log \det$ 运算。当我们试图扩展到更大的数据集时，这很快就会在计算上变得令人望而却步。与其最小化每组增强样本 $\tilde{\vX}_{i}$ 的精确编码率，我们可以通过最大化增强样本特征之间的{\em 余弦相似度}来近似压缩每组特征的效果：\footnote{我们留给读者一个练习，去验证在什么条件下最小化每个类别的编码率 $R^c(\vZ)$ （近似）等价于最大化特征 $\vZ$ 之间的余弦相似度。}

\begin{equation}
    \begin{aligned}
        \max_{\theta } & \;  R_{\epsilon}(\vZ) +
        \frac{\gamma}{N}\sum_{k=1}^N \sum_{i \neq j}
        \vz_{k,j}^\top \vz_{k,i}, \\
        \mbox{s.t.} &  \ \ \, \vz_{i,j} = f(\tau_j(\vx_i),\theta),\
        \|\vz_{i,j}\|_2^2                                          = 1,
    \end{aligned}
    \label{eqn:maximal-rate-reduction-unsupervised-simdino}
\end{equation}
其中 $\gamma > 0$ 是一个超参数，用于控制余弦相似度项的强度，因为它在数量级上与编码率不同。这个目标避免了昂贵的 $\log \det$ 运算，并且具有更好的可扩展性。事实上，这个目标与流行的（无监督）{\em 对比学习}方法密切相关，例如 SimCLR \cite{chen2020simple} 和 DINO \cite{caron2021emerging}，这些方法经过工程设计，能够从海量图像中学习高质量的表征。

正如我们稍后将在第 \ref{ch:applications} 章的第 \ref{sub:contrastive_learning_objective} 节中展示的那样，这个目标 \eqref{eqn:maximal-rate-reduction-unsupervised-simdino} 可以非常高效地实现，以学习高度线性且具有判别性的高度结构化表征 $\vZ$。
% \yima{Robin，也许可以提供一些实验结果，定性地说明学习到的特征具有判别性和线性？比如线性探测（linear probing）的一些结果或分割的视觉结果？（可以说将实现细节留到第 8 章。）这将有助于证明后面第 6 章中 RAE 的合理性。}
学习到的表征在大型真实世界数据集上达到了最先进的性能，与 DINO 学习到的特征相当甚至更好，但不需要 DINO 所需的繁重工程工作。因此，我们将此目标称为 {\em 简化版 DINO}（Simplified DINO，简称 SimDINO）。作为预览，\Cref{tab:dino_imagenet_linear_probing_teaser} 显示，在使用 \(k\)-最近邻准确率和线性探测（即在预训练模型上训练线性分类器）进行 ImageNet-1k 分类时，SimDINO 均优于 DINO。这表明该目标能够产生高度结构化和线性化的表征。我们还可以将这些表征用于下游任务，例如无监督分割\footnote{具体而言，我们采用了 MaskCut \cite{wang2023cut}，这是一种从冻结的视觉骨干网络中提取特征以进行目标检测和实例分割的有效无监督方法。}，如 \Cref{fig:maskcut_segmentation} 所示。SimDINO 的更多实现细节和实验结果将在第 \ref{ch:applications} 章中介绍。

\begin{table}
    \centering
    \begin{tabular}{@{}lcccc@{}} %
        \toprule
        方法 & 模型 & 轮数 & 20-NN & 线性探测
        \\
        \midrule
        DINO & ViT-B & 100 & 72.9 & 76.3 \\
        SimDINO & ViT-B & 100 & \bf 74.9 & \bf 77.3 \\
        DINO & ViT-L & 100 & -- & -- \\
        SimDINO & ViT-L & 100 & \bf 75.6 & \bf 77.4 \\

        \bottomrule
    \end{tabular}\smallskip
    \caption{\small DINO 和 SimDINO 在 ImageNet-1k 测试集上的\textbf{分类性能}，使用了 \(k\)-最近邻准确率（\(k = 20\)）和线性探测。SimDINO 在准确率和训练稳定性方面均优于 DINO。DINO 在 ViT-L 上的训练不稳定且很快发散，而 SimDINO 的训练则很稳定。}
    \label{tab:dino_imagenet_linear_probing_teaser}
\end{table}

\begin{figure}
    \centering
    \includegraphics[width=\linewidth]{chapters/chapter4/figs/maskcut_demo.pdf}
    \caption{\small 来自 DINO ViT-B/16（第 1 行）、SimDINO ViT-B/16（第 2 行）和 SimDINO ViT-L/16（第 3 行）的\textbf{无监督分割结果可视化}。}
    \label{fig:maskcut_segmentation}
\end{figure}

\section{总结与注记}

\paragraph{核心信息。}
在本章中，我们研究了如何从有限样本中学习分布并获得该分布的{\em 可计算}编码和解码方案背后的基本概念和思想。对于连续且可能是低维的分布，我们论证了这种方案必然是{\em 有损的}——因此失真编码率概念非常重要。由于量化而通过失真引入的噪声，与我们在前一章中研究的噪声起着相关但不同的作用。人们可以认为第 \ref{ch:denoising} 章中考虑的噪声位于分布支撑子流形的“外部”，而本章中考虑的噪声位于支撑的“内部”。通俗地说，我们可以将支撑外部的扩散/去噪过程称为“延拓（continuation）”，将支撑内部的扩散/去噪过程称为“渗流（percolation）”。对它们复合效应的更统一的分析和理解有待未来的研究。

在本章中，我们首次引入了为分布学习{\em 表征}的概念。我们考虑了可以说是最简单的情况，即分类的情况。如果将分类标签 $\y$ 视为数据 $\vx$ 的固定简单（严重有损）编码，我们可以寻求数据 $\vx$ 的一致\footnote{也就是说，我们可以从 $\vz$ 中忠实地恢复 $\vx$，正如我们稍后将在第 \ref{ch:consistent-representations} 章中明确“一致”表征的概念那样。}表征 $\vz = f(\vx)$，使其倾向于最大化该表征与这个简单编码之间的互信息：
\begin{equation}
    \max I(\vz, \y) = H(\vz) - H(\vz \mid \y).
\end{equation}
此外，我们要求这种表征 $\vz$ 不仅与 $\vx$ 一致，而且比 $\vx$ 的原始数据分布更加{\em 结构化}——类别是线性可分的。这将使得学习到的表征在用于识别、分割和补全等后续任务时，比原始数据分布高效得多。引入的最大化编码率降低目标是这些要求的一种自然且可计算的实现。在后面的第 \ref{ch:conditional-inference} 章中，我们将看到当 $\y$ 可能是更一般的有损编码（例如图像描述）或与 $\vx$ 相关的随机变量时，如何将此目标推广到更广泛的背景中。

\section{练习与扩展}

\begin{exercise}
    请证明 $\log\det(\cdot)$ 函数的以下性质。
    \begin{enumerate}
        \item 证明
            \begin{align*}
                f(\vX) = \log\det\left(\vX\right)
            \end{align*}
            是一个凹函数。（{\bf 提示：} 函数 $f(\vx)$ 是凹函数当且仅当一维函数 $f(\vx+t\vh)$ 对所有 $\vx$ 和 $\vh$ 都是凹的。）

        \item 证明
            \begin{align*}
                \log\det(\vI + \vX^\top\vX) = \log\det(\vI + \vX\vX^\top)
            \end{align*}

        \item 设 $\vA \in \R^{n\times n}$ 为正定矩阵。请证明：
            \begin{align}
                \log\det\left(\vA \right) = \sum_{i=1}^n \log(\lambda_i),
            \end{align}
            其中 $\lambda_1,\lambda_2,\dots,\lambda_n$ 是 $\vA$ 的特征值。

        \item  证明对于正定矩阵 $\vA$，$\log\det(\cdot)$ 函数的梯度由下式给出：
            \begin{align*}
                \nabla \log\det(\vA) = \vA^{-1},
            \end{align*}

    \end{enumerate}

\end{exercise}

% \begin{exercise}
%     设 $\bm Z \in \R^{d\times N}$ 并定义
%     \begin{equation}
%     R_{\epsilon}(\vZ) := \frac{1}{2}\log\det\left(\I +
%     \alpha\vZ\vZ^{\top}\right),
%     \end{equation}
%     其中 $\alpha > 0$。
%     \begin{enumerate}
%         \item 请计算上述函数的梯度 $\nabla R_{\epsilon}(\bm Z)$。

%         \item 设 $\bm D \in \mathbb{R}^{d \times n}$ 为任意矩阵。请计算上述函数海森矩阵（Hessian）的双线性形式，即：
%         \begin{align*}
%             \nabla^2 R_{\epsilon}(\bm Z)[\bm D, \bm D] :=
% \left\langle \bm D, \lim_{t \to 0} \frac{\nabla R_{\epsilon}(\bm Z
% + t\bm D) - \nabla R_{\epsilon}(\bm Z)}{t}\right\rangle.
%         \end{align*}
%     \end{enumerate}
% \end{exercise}

\begin{exercise}
    设 $\bm Z = [\bm Z_1, \dots, \bm Z_K] \in \mathbb{R}^{d \times N}$，其中每个 $\bm Z_k \in \mathbb{R}^{d \times N_k}$ 且 $N = \sum_{k=1}^K N_k$。

    \begin{enumerate}
        \item 请证明
            \begin{align*}
                \log\det\left( \bm I + \alpha \bm Z \bm Z^\top \right)
                \le
                \sum_{k=1}^K
                \log\det\left( \bm I + \alpha \bm Z_k \bm Z_k^\top \right),
            \end{align*}
            其中 $\alpha > 0$。

        \item 请推导等号成立的充要条件。
    \end{enumerate}

\end{exercise}

\begin{exercise}
    设 $\vz_1$ 和 $\vz_2$ 为两个单位长度的向量。设 $\vZ = [\vz_1, \vz_2]$。证明以下三个目标是等价的：
    \begin{itemize}
        \item $\max \vz_1^\top \vz_2$。
        \item $\min \|\vz_1 - \vz_2\|_2^2$。
        \item $\max \log \det (\I + \vZ\vZ^\top)$。
    \end{itemize}
    然而，当高维空间中有两个以上的向量时，它们可能不再完全等价。$\log\det$ 更为一般，因为它测量了高维空间中所有向量张成的总体积。
\end{exercise}

% {\color{blue}
% \begin{exercise}\label{exer:gaussian_rate_distortion}
%     请证明信源 $X \sim \mathcal{N}(0,\sigma^2)$ 的失真编码率函数为
%     \begin{align*}
%         \mathcal{R}_{\epsilon}(X) =
%         \begin{cases}
%             \frac{1}{2}\log\left( \frac{\sigma^2}{\epsilon^2}
%             \right),\ & 0 \le \epsilon \le \sigma,\\
%             0,\ & \epsilon > \sigma.
%         \end{cases}
%     \end{align*}
%     {\bf 提示：} 该证明包含两个步骤。首先，我们找到失真编码率函数的下界，即对于所有 $\mathbb{E}[(X-\hat{X})^2] \le \epsilon^2$，有
%     \begin{align*}
%         I(X;\hat{X}) \ge  \frac{1}{2}\log\left(
%         \frac{\sigma^2}{\epsilon^2} \right).
%     \end{align*}
%     其次，我们构造 $p(\hat{X} \mid X)$ 来证明这个下界是可达的。如果 $\epsilon \le \sigma$，我们选择
%     \begin{align*}
%         X = \hat{X} + Z,\ \hat{X} \sim \mathcal{N}(0,\sigma^2 -
%         \epsilon^2),\ Z \sim \mathcal{N}(0,\epsilon^2),
%     \end{align*}
%     其中 $\hat{X}$ 和 $Z$ 是独立的。如果 $\epsilon > \sigma$，我们以概率 1 选择 $\hat{X} = 0$。
% \end{exercise}

\begin{exercise}\label{exer:KKT}
    请推导 \Cref{ex:gaussian_rate_distortion} 中引入的高斯随机向量的失真编码率函数。
    % \textbf{提示：} 证明包含以下步骤：
    % \begin{enumerate}
    %     \item 首先，我们找到失真编码率函数的下界，即对于所有 $\mathbb{E}[\|\bm{x}-\hat{\bm x}\|_2^2] \le \epsilon^2$，有
    %         \begin{align}
    %             I(\bm x; \hat{\bm x}) \ge \sum_{i=1}^d
    %             \max\left\{\frac{1}{2}\log\left(
    %             \frac{\sigma_i}{\lambda_i} \right), 0\right\},
    %         \end{align}
    %         其中 $\lambda_i = \mathbb{E}[(x_i - \hat{x}_i)^2]$。

    %     \item 其次，我们将寻找失真编码率函数的问题简化为以下问题：
    %         \begin{align}
    %             \min_{\lambda_1,\dots,\lambda_d}  \sum_{i=1}^d
    %             \max\left\{\frac{1}{2}\log\left(
    %             \frac{\sigma_i}{\lambda_i} \right), 0 \right\}\quad
    %             \mathrm{s.t.}\ \sum_{i=1}^d \lambda_i = \epsilon^2.
    %         \end{align}

    %         % \item 考虑建立拉格朗日函数
    %         % \[
    %         %     \mathcal{L}(\lambda_1,\dots,\lambda_d,\nu) =
    %         % \frac12 \sum_{i=1}^d \log \frac{\sigma_i^2}{\lambda_i}
    %         %     + \nu \left(\sum_{i=1}^d \lambda_i -
    %         % \epsilon^2\right),
    %         % \]
    %         % 其中 \(\lambda_i\) 是沿每个主方向的重建误差的方差。

    %         % \item 计算偏导数 \(\partial \mathcal{L}/\partial \lambda_i = 0\) 以找到 \(\lambda_i\) 和拉格朗日乘子 \(\nu\) 之间的关系。

    %         % \item 应用约束 \(0 \le \lambda_i \le \sigma_i^2\) 来确定哪些方向是饱和的，哪些方向遵循注水（water-filling）分配。

    %         % \item 使用 \(\sum_{i=1}^d \lambda_i = \epsilon^2\) 求解“水位” \(\lambda\)，并验证 \(\lambda_i = \min(\sigma_i^2, \lambda)\)。
    % \end{enumerate}
\end{exercise}
% }

\end{document}
